% This chapter has been modified on 6-4-05.
%\setcounter{chapter}{4}
\choice{\chapter{Distribusi Penting}} 
{\chapter[Distribusi dan Densitas]{Distribusi dan Densitas Penting}}\label{chp 5} 

\section{Distribusi Penting}\label{sec 5.1} 

Dalam bab ini, kita menguraikan distribusi peluang diskret dan densitas peluang kontinu
yang paling sering muncul dalam analisis percobaan. Kita juga akan menunjukkan cara
menyimulasikan distribusi dan densitas tersebut pada komputer.

\subsection*{Distribusi Seragam Diskret}\index{distribusi seragam}\index{fungsi
distribusi!seragam}

Dalam Bab~\ref{chp 1}, kita telah melihat bahwa dalam banyak kasus kita menganggap semua
hasil suatu percobaan berpeluang sama. Jika $X$ adalah peubah acak yang menyatakan hasil
percobaan semacam ini, kita mengatakan bahwa $X$ terdistribusi secara seragam. Jika ruang
sampel $S$ berukuran $n$, dengan $0 < n < \infty$, fungsi distribusi $m(\omega)$
didefinisikan sebagai $1/n$ untuk setiap $\omega \in S$. Seperti semua distribusi peluang
diskret yang dibahas dalam bab ini, percobaan tersebut dapat disimulasikan pada komputer
dengan program {\bf GeneralSimulation}. Namun, dalam kasus ini kita dapat menggunakan
algoritma yang lebih cepat. (Algoritma ini telah diuraikan dalam Bab~\ref{chp 1}; uraian
tersebut kita ulang di sini demi kelengkapan.) Ekspresi
$$1 + \lfloor n\,(rnd)\rfloor$$ 
bernilai setiap bilangan bulat antara 1 dan $n$ dengan peluang $1/n$ (notasi
$\lfloor x \rfloor$ menyatakan bilangan bulat terbesar yang tidak melebihi $x$). Jadi,
jika hasil-hasil yang mungkin dari percobaan diberi label
$\omega_1\ \omega_2,\ \ldots,\ \omega_n$, kita menggunakan ekspresi di atas untuk
menyatakan subskrip hasil percobaan.
\par Jika ruang sampelnya merupakan himpunan tak hingga terhitung, seperti himpunan
bilangan bulat positif, tidak mungkin ada percobaan yang seragam pada himpunan tersebut
(lihat Latihan~\ref{exer 5.1.102}). Jika ruang sampelnya merupakan himpunan tak terhitung
dengan panjang positif dan terhingga, seperti interval $[0, 1]$, kita menggunakan fungsi
densitas kontinu (lihat Bagian~\ref{sec 5.2}).

\subsection*{Distribusi Binomial}\index{fungsi distribusi!binomial}\index{distribusi
binomial}

Distribusi binomial dengan parameter $n$, $p$, dan $k$ telah didefinisikan dalam
Bab~\ref{chp 3}. Distribusi ini merupakan distribusi peubah acak yang mencacah banyaknya
kepala yang muncul ketika sebuah koin dilempar $n$ kali, dengan anggapan bahwa pada setiap
lemparan peluang munculnya kepala adalah $p$. Fungsi distribusinya diberikan oleh rumus
$$b(n, p, k) = {n \choose k}p^k q^{n-k}\ ,$$ dengan $q = 1 - p$.
\par Salah satu cara langsung untuk menyimulasikan peubah acak binomial $X$ adalah
menghitung jumlah $n$ peubah acak $0-1$ yang saling bebas, yang masing-masing bernilai 1
dengan peluang $p$. \choice{}{ Cara ini memerlukan $n$ pemanggilan generator bilangan acak
untuk memperoleh satu nilai peubah acak. Ketika $n$ relatif besar (misalnya sedikitnya
30), Teorema Limit Pusat (lihat Bab~\ref{chp 9}) menyatakan bahwa distribusi binomial
didekati dengan baik oleh fungsi densitas normal yang bersesuaian (yang didefinisikan
dalam Bagian~\ref{sec 5.2}) dengan parameter $\mu = np$ dan
$\sigma = \sqrt{npq}$. Jadi, dalam kasus ini kita dapat menghitung suatu nilai $Y$ dari
peubah acak normal dengan parameter tersebut, dan jika $-1/2 \le Y < n+1/2$, kita dapat
menggunakan nilai
$$\lfloor Y + 1/2 \rfloor$$ untuk menyatakan peubah acak $X$. Jika $Y < -1/2$ atau $Y
> n + 1/2$, kita menolak $Y$ dan menghitung nilai lain. Dalam bagian berikutnya kita akan
melihat cara menyimulasikan peubah acak normal dengan cepat.}


\subsection*{Distribusi Geometrik}\index{distribusi geometrik}\index{fungsi
distribusi!geometrik}

Tinjaulah suatu proses percobaan Bernoulli yang dilanjutkan hingga tak hingga banyaknya
percobaan; sebagai contoh, sebuah koin yang dilempar dalam barisan tak hingga.
\choice{Ukuran peluang untuk suatu proses yang memerlukan tak hingga banyaknya percobaan
ditentukan berdasarkan peluang yang hanya memerlukan sejumlah terhingga
percobaan\footnote{Lihat Bagian 2.2 dalam buku lengkap Grinstead--Snell}}{Dalam
Bagian~\ref{sec 2.2} telah kita tunjukkan cara menetapkan distribusi peluang pada pohon
tak hingga.} Dengan demikian, kita dapat menentukan distribusi sebarang peubah acak $X$
yang berkaitan dengan percobaan tersebut, asalkan $P(X = a)$ dapat dihitung berdasarkan
sejumlah terhingga percobaan. Sebagai contoh, misalkan $T$ adalah banyaknya percobaan
sampai dengan dan termasuk keberhasilan pertama. Maka
\begin{eqnarray*} P(T = 1) & = & p\ , \\ P(T = 2) & = & qp\ , \\ P(T = 3) & = & q^2p\ , \\
\end{eqnarray*}
dan secara umum,
$$P(T = n) = q^{n-1}p\ .$$ 
Untuk menunjukkan bahwa ini merupakan suatu distribusi, kita harus menunjukkan bahwa
$$ p + qp + q^2p + \cdots = 1\ .
$$ 
Ekspresi di ruas kiri hanyalah deret geometrik dengan suku pertama $p$ dan rasio tetap
$q$, sehingga jumlahnya adalah
$${p\over{1-q}}$$ yang sama dengan 1.
\par Dalam Gambar~\ref{fig 5.4} kita telah menggambar distribusi ini menggunakan program
{\bf GeometricPlot} untuk kasus $p = .5$ dan $p = .2$. Terlihat bahwa ketika $p$
mengecil, kita semakin mungkin memperoleh nilai $T$ yang besar, sebagaimana dapat diduga.
Dalam kedua kasus, nilai $T$ yang paling mungkin adalah 1. Hal ini selalu berlaku karena
$$
\frac {P(T = j + 1)}{P(T = j)} = q < 1\ .
$$
\putfig{4.75truein}{PSfig5-4}{Distribusi geometrik.}{fig 5.4} 
\par
Secara umum, jika $0 < p < 1$ dan $q = 1 - p$, kita mengatakan bahwa peubah acak $T$
memiliki \emx {distribusi geometrik} jika
$$P(T = j) = q^{j - 1}p\ ,$$ untuk $j = 1,\ 2,\ 3,\ \ldots$ .

\par Untuk menyimulasikan distribusi geometrik dengan parameter $p$, kita cukup
menghasilkan suatu barisan bilangan acak dalam $[0, 1)$ dan berhenti ketika salah satu
nilainya tidak melebihi $p$. Namun, untuk nilai $p$ yang kecil, cara ini memakan waktu
(rata-rata memerlukan $1/p$ langkah). Sekarang kita uraikan suatu metode yang waktu
jalannya tidak bergantung pada besar $p$. Definisikan $Y$ sebagai bilangan bulat terkecil
yang memenuhi pertidaksamaan
\begin{equation}
1 - q^Y \ge rnd\ .\label{eq 5.3}
\end{equation} 
Maka kita memiliki
\begin{eqnarray*} P(Y = j) & = & P\Bigl(1 - q^j \ge rnd > 1 - q^{j-1}\Bigr) \\ & = &
q^{j-1} - q^j \\ & = & q^{j-1}(1-q) \\ & = & q^{j-1}p\ . \\
\end{eqnarray*} Jadi, $Y$ berdistribusi geometrik dengan parameter $p$. Untuk
menghasilkan $Y$, kita hanya perlu menyelesaikan Persamaan~\ref{eq 5.3} terhadap $Y$.
Kita memperoleh
$$Y = \Biggl\lceil {{\log(1-rnd)}\over{\log\ q}}\Biggr\rceil\ ,$$
dengan notasi $\lceil x \rceil$ yang menyatakan bilangan bulat terkecil yang lebih besar
dari atau sama dengan $x$. Karena $\log(1-rnd)$ dan $\log(rnd)$ berdistribusi identik,
$Y$ juga dapat dihasilkan menggunakan persamaan
$$Y = \Biggl\lceil {{\log\ rnd}\over{\log\ q}}\Biggr\rceil\ .$$

\begin{example}\label{exam 5.1}
Distribusi geometrik berperan penting dalam teori antrean\index{antrean}. Sebagai contoh,
misalkan sebaris pelanggan menunggu layanan di sebuah loket. Sering kali dianggap bahwa
dalam setiap satuan waktu yang singkat, terdapat 0 atau 1 pelanggan baru yang tiba di
loket. Peluang bahwa seorang pelanggan tiba adalah $p$, sedangkan peluang bahwa tidak ada
pelanggan yang tiba adalah $q = 1 - p$. Dengan demikian, waktu $T$ hingga kedatangan
berikutnya memiliki distribusi geometrik. Wajar jika kita menanyakan peluang bahwa tidak
ada pelanggan yang tiba dalam $k$ satuan waktu berikutnya, yakni $P(T > k)$. Peluang ini
diberikan oleh
\begin{eqnarray*} P(T > k) = \sum_{j = k+1}^\infty q^{j-1}p & = & q^k(p + qp + q^2p + \cdots)
\\
                                    & = & q^k\ .
\end{eqnarray*} Peluang ini juga dapat ditemukan dengan memperhatikan bahwa kita meminta
tidak adanya keberhasilan (yakni, kedatangan) dalam suatu barisan $k$ satuan waktu
berturut-turut, dengan peluang keberhasilan pada setiap satuan waktu sebesar $p$. Jadi,
peluangnya tepat $q^k$, karena kedatangan pada dua satuan waktu mana pun merupakan
kejadian yang saling bebas.
\par
Sering kali dianggap bahwa lama waktu yang diperlukan untuk melayani seorang pelanggan
juga memiliki distribusi geometrik, tetapi dengan nilai $p$ yang berbeda. Hal ini
menyiratkan suatu sifat waktu layanan yang cukup khusus. Untuk melihatnya, mari kita
hitung peluang bersyarat
$$ P(T > r + s\,|\,T > r) = \frac{P(T > r + s)}{P(T > r)} = \frac {q^{r + s}}{q^r} =
q^s\ .
$$ 
Jadi, peluang bahwa layanan pelanggan tersebut memerlukan tambahan $s$ satuan waktu tidak
bergantung pada lama waktu $r$ selama pelanggan itu telah dilayani. Karena penafsiran ini,
sifat tersebut disebut ``sifat tanpa ingatan," dan distribusi eksponensial juga
memilikinya. (Untungnya, tidak banyak tempat pelayanan yang memiliki sifat ini.)
\end{example}

\subsection*{Distribusi Binomial Negatif}\index{distribusi binomial
negatif}\index{fungsi distribusi!binomial negatif}

Misalkan kita mempunyai sebuah koin yang ketika dilempar menghasilkan kepala dengan
peluang $p$. Kita tetapkan suatu bilangan bulat positif $k$, lalu melempar koin hingga
kepala ke-$k$ muncul. Misalkan $X$ menyatakan banyaknya lemparan. Ketika $k = 1$, $X$
berdistribusi geometrik. Untuk $k$ umum, kita mengatakan bahwa $X$ memiliki distribusi
binomial negatif. Sekarang kita hitung distribusi peluang $X$. Jika $X = x$, tepat
$k-1$ kepala harus muncul dalam $x-1$ lemparan pertama, dan kepala harus muncul pada
lemparan ke-$x$. Terdapat
$${{x-1}\choose{k-1}}$$ 
barisan sepanjang $x$ dengan sifat-sifat ini, dan setiap barisan tersebut diberi peluang
yang sama, yakni
$$p^{k-1}q^{x-k}\ .$$ Oleh karena itu, jika kita definisikan
$$u(x, k, p) = P(X = x)\ ,$$
maka
$$u(x, k, p) = {{x-1}\choose{k-1}}p^kq^{x-k}\ .$$
\par Kita dapat menyimulasikannya pada komputer dengan menyimulasikan pelemparan koin.
Namun, algoritma berikut umumnya jauh lebih cepat. Perhatikan bahwa $X$ dapat dipahami
sebagai jumlah $k$ hasil percobaan yang masing-masing berdistribusi geometrik dengan
parameter $p$. Jadi, kita dapat menggunakan jumlah berikut untuk menghasilkan $X$:
$$\sum_{j = 1}^k \Biggl\lceil {{\log\ rnd_j}\over{\log\ q}}\Biggr\rceil\ .$$

\begin{example} Sebuah koin adil dilempar hingga kepala muncul untuk kedua kalinya.
Distribusi banyaknya lemparan adalah $u(x, 2, p)$. Jadi, peluang bahwa diperlukan $x$
lemparan untuk memperoleh dua kepala ditemukan dengan menetapkan $k = 2$ dalam rumus di
atas. Kita memperoleh
$$  u(x, 2, 1/2) =  {{x-1} \choose 1} \frac 1{2^x}\ ,$$ 
untuk $x = 2, 3, \ldots\ $.
\par
Dalam Gambar~\ref{fig 7.2} ditampilkan grafik distribusi untuk~$k = 2$ dan~$p = .25$.
Perhatikan bahwa distribusi tersebut cukup asimetris, dengan ekor panjang yang
mencerminkan kemungkinan munculnya nilai~$x$ yang besar.
\end{example}

\putfig{4.5truein}{PSfig7-2-1}{Distribusi binomial negatif dengan $k = 2$ dan $p = .25$.}{fig 7.2} 


\subsection*{Distribusi Poisson}\index{distribusi Poisson}\index{fungsi
distribusi!Poisson}

Distribusi Poisson muncul dalam banyak keadaan. Dapat dikatakan bahwa distribusi ini
merupakan salah satu dari tiga distribusi peluang diskret terpenting (dua lainnya adalah
distribusi seragam dan binomial). Distribusi Poisson dapat dipandang sebagai distribusi
yang timbul dari distribusi binomial atau dari densitas eksponensial. Sekarang kita akan
menjelaskan hubungannya dengan yang pertama; hubungannya dengan yang kedua akan dijelaskan
dalam bagian berikutnya.
\par Misalkan kita menghadapi suatu keadaan ketika jenis kejadian tertentu berlangsung
secara acak selama suatu selang waktu. Sebagai contoh, kejadian yang kita minati mungkin
berupa panggilan telepon yang masuk ke kantor polisi di sebuah kota besar. Kita ingin
memodelkan keadaan ini agar dapat meninjau peluang kejadian seperti munculnya lebih dari
10 panggilan telepon dalam selang 5 menit. Dalam contoh kita, agaknya terdapat lebih
banyak panggilan masuk antara pukul 6:00 dan 7:00~{\footnotesize malam} daripada antara
pukul 4:00 dan 5:00~{\footnotesize pagi}, dan fakta ini tentu memengaruhi peluang di atas.
Jadi, agar peluang semacam itu dapat dihitung, kita harus menganggap bahwa laju rata-rata,
yakni rata-rata banyaknya kejadian per menit, merupakan suatu konstanta. Laju ini kita
nyatakan dengan $\lambda$. (Jadi, dalam suatu selang waktu 5 menit, kita mengharapkan
sekitar $5\lambda$ kejadian.) Artinya, jika model kita diterapkan pada kedua rentang waktu
di atas, kita cukup menggunakan laju yang berbeda untuk masing-masing rentang sehingga
memperoleh dua peluang yang berbeda bagi kejadian tersebut.
\par Anggapan kita berikutnya adalah bahwa banyaknya kejadian dalam dua selang waktu yang
tidak tumpang tindih saling bebas. Dalam contoh kita, ini berarti bahwa kejadian terdapat
$j$ panggilan antara pukul 5:00 dan 5:15~{\footnotesize malam} serta kejadian terdapat $k$
panggilan antara pukul 6:00 dan 6:15~{\footnotesize malam} pada hari yang sama saling
bebas.
\par Kita dapat menggunakan distribusi binomial untuk memodelkan keadaan ini. Bayangkan
suatu selang waktu tertentu dibagi menjadi $n$ subselang yang sama panjang. Jika
subselang tersebut cukup singkat, kita dapat menganggap bahwa peluang terjadinya dua atau
lebih kejadian dalam satu subselang dapat diabaikan dibandingkan dengan peluang terjadinya
paling banyak satu kejadian. Jadi, pada setiap subselang kita menganggap terdapat 0 atau 1
kejadian. Dengan demikian, barisan subselang dapat dipandang sebagai barisan percobaan
Bernoulli, dengan keberhasilan bersesuaian dengan suatu kejadian dalam subselang tersebut.
\par Untuk menentukan nilai $p$ yang tepat, yakni peluang suatu kejadian dalam satu
subselang, kita bernalar sebagai berikut. Rata-rata terdapat $\lambda t$ kejadian dalam
selang waktu sepanjang $t$. Jika selang waktu ini dibagi menjadi $n$ subselang, berdasarkan
penafsiran percobaan Bernoulli kita mengharapkan terdapat $np$ kejadian. Jadi, kita ingin
$$\lambda t = n p\ ,$$ sehingga
$$p = {{\lambda t}\over{n}}\ .$$
\par Sekarang kita hendak meninjau peubah acak $X$ yang mencacah banyaknya kejadian dalam
suatu selang waktu tertentu. Kita ingin menghitung distribusi $X$. Untuk memudahkan
perhitungan, kita anggap selang waktunya memiliki panjang 1; untuk selang waktu dengan
panjang sembarang $t$, lihat Latihan~\ref{exer 5.1.26}. Kita mengetahui bahwa
$$P(X = 0) = b(n, p, 0) = (1 - p)^n = \Bigl(1 - {\lambda \over n}\Bigr)^n\ .$$ Untuk
nilai $n$ yang besar, nilai ini kira-kira $e^{-\lambda}$. Mudah dihitung bahwa untuk
setiap $k$ tetap, kita memiliki
$${{b(n, p, k)}\over{b(n, p, k-1)}} = {{\lambda - (k-1)p}\over{kq}}$$ yang, untuk $n$
besar (dan dengan demikian $p$ kecil), kira-kira sama dengan $\lambda/k$. Jadi, kita
memiliki
$$P(X = 1) \approx \lambda e^{-\lambda}\ ,$$ dan secara umum,
\begin{equation} P(X = k) \approx {{\lambda^k}\over{k!}} e^{-\lambda}\ .\label{eq 5.1}
\end{equation} Distribusi di atas adalah distribusi Poisson. Perlu diperiksa bahwa
distribusi yang diberikan dalam Persamaan~\ref{eq 5.1} benar-benar \emx {merupakan}
suatu distribusi, yakni bahwa nilai-nilainya non-negatif dan berjumlah 1. (Lihat
Latihan~\ref{exer 5.1.27}.)
\par Distribusi Poisson digunakan sebagai aproksimasi distribusi binomial ketika
parameter $n$ dan $p$ masing-masing besar dan kecil (lihat Contoh~\ref{exam 5.3} dan
\ref{exam 5.5}). Namun, distribusi Poisson juga muncul dalam keadaan ketika parameter
$n$ dan $p$ mungkin tidak mudah ditafsirkan atau diukur (lihat Contoh~\ref{exam 5.5.5}).

\begin{example}\label{exam 5.3} Seorang penata huruf\index{penata huruf} rata-rata membuat
satu kesalahan per 1000 kata. Anggaplah ia sedang menata sebuah buku dengan 100 kata per
halaman. Misalkan $S_{100}$ adalah banyaknya kesalahan yang dibuatnya pada satu halaman.
Distribusi peluang eksak bagi~$S_{100}$ kemudian diperoleh dengan memandang $S_{100}$
sebagai hasil 100 percobaan Bernoulli dengan $p = 1/1000$. Nilai harapan~$S_{100}$ adalah
$\lambda = 100(1/1000) = .1$. Peluang eksak bahwa $S_{100} = j$ adalah
$b(100,1/1000,j)$, dan aproksimasi Poisson\index{aproksimasi Poisson terhadap distribusi\\
binomial} adalah
$$
\frac {e^{-.1}(.1)^j}{j!}.
$$ Dalam Tabel~\ref{table 5.1}, untuk berbagai nilai $n$ dan $p$, diberikan nilai eksak
yang dihitung menggunakan distribusi binomial dan aproksimasi Poisson.
\begin{table}
\centering
\begin{tabular}{|l|c|c|c|c|c|c|}
\hline
&Poisson&Binomial&Poisson&Binomial&Poisson&Binomial \\
&&$n = 100$&&$n = 100$&&$n = 1000$\\
$j$&$\lambda = .1$&$p = .001$&$\lambda = 1$&$p = .01$&$\lambda = 10$&$p = .01$\\
\hline
0&.9048&.9048&.3679&.3660&.0000&.0000\\
1&.0905&.0905&.3679&.3697&.0005&.0004\\
2&.0045&.0045&.1839&.1849&.0023&.0022\\
3&.0002&.0002&.0613&.0610&.0076&.0074\\
4&.0000&.0000&.0153&.0149&.0189&.0186\\
5&&&.0031&.0029&.0378&.0374\\
6&&&.0005&.0005&.0631&.0627\\
7&&&.0001&.0001&.0901&.0900\\
8&&&.0000&.0000&.1126&.1128\\
9&&&&&.1251&.1256\\
10&&&&&.1251&.1257\\
11&&&&&.1137&.1143\\
12&&&&&.0948&.0952\\
13&&&&&.0729&.0731\\
14&&&&&.0521&.0520\\
15&&&&&.0347&.0345\\
16&&&&&.0217&.0215\\
17&&&&&.0128&.0126\\
18&&&&&.0071&.0069\\
19&&&&&.0037&.0036\\
20&&&&&.0019&.0018\\
21&&&&&.0009&.0009\\
22&&&&&.0004&.0004\\
23&&&&&.0002&.0002\\
24&&&&&.0001&.0001\\
25&&&&&.0000&.0000\\
\hline
\end{tabular}
\caption{Aproksimasi Poisson terhadap distribusi binomial.}
\label{table 5.1}
\end{table}
\end{example}

\begin{example}\label{exam 5.5} Dalam bukunya,\footnote{ibid., p.~161.}
Feller\index{FELLER, W.} membahas statistik jatuhnya bom terbang\index{bom terbang} di
London bagian selatan selama Perang Dunia Kedua.
\par
Anggaplah Anda tinggal di sebuah kawasan berukuran 10~blok kali 10~blok sehingga seluruh
kawasan terbagi menjadi 100 persegi kecil. Berapakah peluang bahwa persegi tempat Anda
tinggal tidak terkena bom jika seluruh kawasan dihantam 400 bom?
\par
Kita menganggap bahwa sebuah bom tertentu akan menghantam persegi Anda dengan
peluang~1/100. Karena terdapat 400 bom, banyaknya hantaman yang diterima persegi Anda
dapat dipandang sebagai banyaknya \emx {keberhasilan} dalam suatu proses percobaan
Bernoulli dengan $n = 400$ dan $p = 1/100$. Jadi, kita dapat menggunakan distribusi
Poisson dengan $\lambda = 400 \cdot 1/100 = 4$ untuk mengaproksimasi peluang bahwa
persegi Anda menerima $j$~hantaman. Peluang ini adalah $p(j) = e^{-4} 4^j/j!$.
Banyaknya persegi yang diharapkan menerima tepat $j$~hantaman kemudian adalah
$100 \cdot p(j)$. Mudah untuk menulis program {\bf LondonBombs} yang menyimulasikan
keadaan ini dan membandingkan banyaknya persegi dengan $j$~hantaman yang diharapkan
dengan banyaknya yang diamati. Dalam Latihan~\ref{exer 9.2.15}, Anda diminta
membandingkan data yang benar-benar diamati dengan data yang diprediksi oleh distribusi
Poisson.
\par
Dalam Gambar~\ref{fig 5.1.5}, ditampilkan hantaman hasil simulasi beserta grafik batang
runcing yang memperlihatkan frekuensi teramati dan terprediksi. Frekuensi teramati
ditampilkan sebagai persegi, sedangkan frekuensi terprediksi ditampilkan sebagai titik.
\putfig{3.5truein}{PSfig5-1-5}{Jatuhnya bom terbang.}{fig 5.1.5} 
\end{example} 
Jika pembaca tidak ingin meninjau bom terbang, sebagai gantinya dapat dipertimbangkan
keadaan serupa yang melibatkan kue kering dan kismis. Anggaplah kita telah membuat adonan
yang cukup untuk 500 kue. Kita memasukkan 600 kismis ke dalam adonan dan mengaduknya
hingga rata. Salah satu cara memandang keadaan ini adalah bahwa kita mempunyai 500 kue,
lalu setelah menata kue-kue tersebut dalam kisi di atas meja, kita melemparkan 600
kismis ke arah kue-kue itu. (Lihat Latihan~\ref{exer 5.1.29}.)
\begin{example}\label{exam 5.5.5} Misalkan dalam suatu volume darah tetap $A$, rata-rata
orang memiliki 40 sel darah putih. Misalkan $X$ adalah peubah acak yang memberikan
banyaknya sel darah putih dalam sampel acak berukuran $A$ dari seseorang yang dipilih
secara acak. Kita dapat memandang $X$ sebagai peubah berdistribusi binomial, dengan setiap
sel darah putih dalam tubuh menyatakan satu percobaan. Jika suatu sel darah putih muncul
dalam sampel, percobaan yang bersesuaian dengan sel darah tersebut merupakan suatu
keberhasilan. Nilai $p$ kemudian harus diambil sebagai rasio $A$ terhadap volume darah
total orang tersebut, dan $n$ adalah banyaknya sel darah putih dalam tubuhnya. Tentu saja,
dalam praktiknya kedua parameter ini tidak mudah diukur secara akurat, tetapi bilangan 40
kiranya mudah diukur. Untuk rata-rata orang, kita kemudian memiliki $40 = np$, sehingga
$X$ dapat dipandang berdistribusi Poisson dengan parameter $\lambda = 40$. Dalam kasus
ini, keadaan tersebut lebih mudah dimodelkan menggunakan distribusi Poisson daripada
distribusi binomial.
\end{example}
\par Cara yang baik untuk menyimulasikan peubah acak Poisson pada komputer adalah
memanfaatkan hubungan antara distribusi Poisson dan densitas eksponensial. Hubungan ini
beserta algoritma simulasi yang dihasilkannya akan diuraikan dalam bagian berikutnya.

\subsection*{Distribusi Hipergeometrik}\index{distribusi hipergeometrik}
\index{fungsi distribusi!hipergeometrik}

Misalkan kita mempunyai sekumpulan $N$ bola, dengan $k$ bola merah dan $N-k$ bola
biru. Kita memilih $n$ bola tanpa pengembalian, lalu mendefinisikan $X$ sebagai
banyaknya bola merah dalam sampel. Distribusi $X$ disebut distribusi hipergeometrik.
Distribusi ini bergantung pada tiga parameter, yaitu $N$, $k$, dan $n$. Tampaknya
tidak ada notasi baku untuk distribusi ini; kita akan memakai notasi $h(N, k, n, x)$
untuk menyatakan $P(X = x)$. Peluang ini dapat dicari dengan memperhatikan bahwa
terdapat
$${N \choose n}$$ sampel berbeda yang berukuran $n$. Banyaknya sampel yang memuat
tepat $x$ bola merah diperoleh dengan mengalikan banyaknya cara memilih $x$ bola
merah dari $k$ bola merah dengan banyaknya cara memilih $n-x$ bola biru dari $N-k$
bola biru. Jadi,
$$h(N, k, n, x) = {{{{k}\choose{x}}{{N-k}\choose{n-x}}}\over{{{N}\choose{n}}}}\ .$$ 
Distribusi ini dapat diperumum untuk keadaan dengan lebih dari dua jenis objek.
(Lihat Latihan~\ref{exer 5.1.24}.)
\par 
Jika $N$ dan $k$ menuju $\infty$ sedemikian rupa sehingga rasio $k/N$ tetap, maka
distribusi hipergeometrik menuju distribusi binomial dengan parameter $n$ dan
$p = k/N$. Hal ini masuk akal karena, jika $N$ dan $k$ jauh lebih besar daripada
$n$, pemilihan sampel dengan atau tanpa pengembalian seharusnya tidak banyak
memengaruhi peluangnya, sedangkan percobaan pemilihan dengan pengembalian menghasilkan
peubah acak berdistribusi binomial (lihat Latihan~\ref{exer 5.1.124}).
\par Contoh penggunaan distribusi ini diberikan dalam Latihan~\ref{exer
5.1.21} dan \ref{exer 5.1.22}. Berikut ini kita berikan contoh lain yang melibatkan distribusi
hipergeometrik. Contoh tersebut menggambarkan suatu uji statistika yang disebut uji
eksak Fisher\index{uji eksak Fisher}.
\begin{example}\label{exam 5.6} Sering kali kita ingin meninjau dua ciri, misalnya
warna mata dan warna rambut, lalu menanyakan apakah terdapat hubungan di antara
keduanya. Dua ciri berhubungan jika pengetahuan tentang nilai salah satu ciri pada
seseorang memungkinkan kita memprediksi nilai ciri lainnya pada orang tersebut.
Semakin kuat hubungannya, semakin akurat prediksinya. Jika tidak terdapat hubungan
di antara kedua ciri, kita mengatakan bahwa keduanya saling bebas. Dalam contoh ini,
kita akan memakai ciri jenis kelamin dan partai politik. Kita akan mengasumsikan bahwa
hanya ada dua kemungkinan jenis kelamin, perempuan dan laki-laki, serta hanya dua
kemungkinan partai politik, Demokrat dan Republikan.
\par Misalkan kita telah mengumpulkan data mengenai ciri-ciri ini. Untuk menguji
apakah terdapat hubungan di antara keduanya, mula-mula kita mengasumsikan bahwa
keduanya tidak berhubungan. Asumsi ini menghasilkan suatu himpunan data ``harapan'',
yang pengetahuan tentang nilai salah satu cirinya tidak membantu memprediksi nilai
ciri lainnya. Himpunan data yang kita kumpulkan biasanya berbeda dari himpunan data
harapan ini. Jika perbedaannya cukup besar, kita cenderung menolak asumsi bahwa kedua
ciri itu saling bebas. Untuk memperjelas makna ``cukup besar'', kita menentukan
himpunan-himpunan data yang berbeda dari himpunan data harapan sekurang-kurangnya
sebesar perbedaan data kita, lalu menghitung peluang bahwa salah satu himpunan data
tersebut terjadi di bawah asumsi bahwa kedua ciri saling bebas. Jika peluang ini
kecil, kecil pula kemungkinannya bahwa perbedaan antara himpunan data yang kita
kumpulkan dan himpunan data harapan sepenuhnya disebabkan oleh kebetulan.
\par Misalkan kita telah mengumpulkan data yang ditampilkan pada Tabel~\ref{table 5.2}.
\begin{table}
\centering
\begin{tabular}{|l|c|c|c|}
\hline
       & \hspace{.15in}Demokrat\hspace{.15in} & \hspace{.1in}Republikan\hspace{.1in} & \\ \hline 
\hspace{.1in}Perempuan & 24                  &\hspace{.02in} 4 & \hspace{.2in} 28  \\ \hline 
\hspace{.1in}Laki-laki   &\hspace{.075in}8     & 14               & \hspace{.2in} 22  \\ \hline
       & 32                  & 18               & \hspace{.2in} 50  \\ \hline
\end{tabular}
\caption{Data teramati.}
\label{table 5.2}
\end{table}
Jumlah pada setiap baris dan kolom disebut total marginal, atau marjinal. Selanjutnya,
kita akan menyatakan jumlah baris dengan $t_{11}$ dan $t_{12}$, serta jumlah kolom
dengan $t_{21}$ dan $t_{22}$. Entri ke-$ij$ pada tabel akan dinyatakan dengan
$s_{ij}$. Terakhir, ukuran himpunan data akan dinyatakan dengan $n$. Jadi, tabel data
umum berbentuk seperti yang ditampilkan pada Tabel~\ref{table 5.3}.
\begin{table}
\centering
\begin{tabular}{|l|c|c|c|}
\hline
 & \hspace{.15in}Demokrat\hspace{.15in} & \hspace{.1in}Republikan\hspace{.1in} & \\ \hline 
\hspace{.1in}Perempuan & $s_{11}$ & $s_{12}$ & \hspace{.2in}$t_{11}$\hspace{.2in}\\ \hline
\hspace{.1in}Laki-laki   & $s_{21}$ & $s_{22}$ & \hspace{.2in}$t_{12}$\\ \hline
                    & $t_{21}$ & $t_{22}$ & \hspace{.2in}$n$\\ \hline
\end{tabular}
\caption{Tabel data umum.}
\label{table 5.3}
\end{table}
Sekarang kita jelaskan model yang akan digunakan untuk membangun himpunan data
``harapan''. Dalam model ini, kita mengasumsikan bahwa kedua ciri saling bebas. Kita
kemudian memasukkan $t_{21}$ bola kuning dan $t_{22}$ bola hijau, yang masing-masing
bersesuaian dengan marjinal Demokrat dan Republikan, ke dalam sebuah guci. Kita menarik
$t_{11}$ bola tanpa pengembalian dari guci dan menyebut bola-bola ini perempuan.
Sebanyak $t_{12}$ bola yang tersisa di dalam guci disebut laki-laki. Dalam kasus
khusus yang sedang ditinjau, peluang memperoleh data aktual menurut model ini
diberikan oleh
$${{{{32}\choose{24}}{{18}\choose{4}}}\over{{{50}\choose{28}}}}\ ,$$
yakni suatu nilai distribusi hipergeometrik.
\par
Sekarang kita siap membangun himpunan data harapan. Jika kita memilih 28 dari 50
bola, secara rata-rata kita mengharapkan persentase bola kuning dalam sampel sama
dengan persentasenya dalam guci. Jadi, secara rata-rata kita mengharapkan
$28(32/50) = 17.92 \approx 18$ bola kuning dalam sampel. (Lihat
Latihan~\ref{exer 6.1.37}.) Nilai-nilai harapan lainnya dihitung dengan cara yang
persis sama. Dengan demikian, himpunan data harapan ditampilkan pada
Tabel~\ref{table 5.4}.
\begin{table}
\centering
\begin{tabular}{|l|c|c|c|}
\hline
 & \hspace{.15in}Demokrat\hspace{.15in} & \hspace{.1in}Republikan \hspace{.1in} & \\ \hline 
\hspace{.1in}Perempuan & 18 & 10           & \hspace{.2in}28\hspace{.2in}\\ \hline 
\hspace{.1in}Laki-laki   & 14 & \hspace{.075in}8  & \hspace{.2in}22 \\ \hline
 & 32 & 18 & \hspace{.2in}50\\ \hline
\end{tabular}
\caption{Data harapan.}
\label{table 5.4}
\end{table}
Perhatikan bahwa nilai $s_{11}$ menentukan tiga nilai lainnya dalam tabel karena
semua marjinal tetap. Jadi, ketika meninjau himpunan-himpunan data yang mungkin muncul
dalam model ini, kita cukup meninjau berbagai kemungkinan nilai $s_{11}$. Dalam kasus
khusus ini, berapakah peluang menarik tepat $a$ bola kuning, yaitu berapakah peluang
bahwa $s_{11} = a$? Peluangnya adalah
\begin{equation}{{{{32}\choose{a}}{{18}\choose{28-a}}}\over{{{50}\choose{28}}}}\ .
\label{eq 5.65}
\end{equation}
\par Sekarang kita siap memutuskan apakah data aktual kita berbeda dari himpunan data
harapan dengan selisih yang lebih besar daripada yang secara wajar dapat dikaitkan
dengan kebetulan semata. Banyaknya perempuan Demokrat yang diharapkan adalah 18,
sedangkan banyaknya dalam data aktual adalah 24. Himpunan data lain yang menyimpang
lebih jauh dari himpunan data harapan bersesuaian dengan keadaan ketika
banyaknya perempuan Demokrat adalah 25, 26, 27, atau 28. Jadi, untuk memperoleh
peluang yang diperlukan, kita menjumlahkan bentuk dalam (\ref{eq 5.65}) dari
$a = 24$ sampai $a = 28$. Kita memperoleh nilai $.000395$. Dengan demikian, kita
harus menolak hipotesis bahwa kedua ciri saling bebas.
\end{example}
\par
Terakhir, kita beralih ke cara menyimulasikan peubah acak hipergeometrik $X$.
Misalkan parameter $X$ adalah $N$, $k$, dan $n$. Bayangkan kita mempunyai sekumpulan
$N$ bola yang diberi label 1 sampai $N$. Kita menetapkan $k$ bola pertama berwarna
merah dan sisanya berwarna biru. Misalkan kita telah memilih $m$ bola dan $j$ di
antaranya berwarna merah. Berarti tersisa $k-j$ bola merah dari total $N-m$ bola.
Dengan demikian, pilihan berikutnya berwarna merah dengan peluang
$${{k-j}\over{N-m}}\ .$$ Pada tahap ini, kita memilih suatu bilangan acak dalam
$[0, 1]$ dan melaporkan bahwa bola merah telah dipilih jika dan hanya jika bilangan
acak itu tidak melebihi bentuk di atas. Kemudian kita memperbarui nilai $m$ dan $j$
dan melanjutkan sampai $n$ bola telah dipilih.

\subsection*{Distribusi Benford}\index{distribusi Benford}\index{fungsi distribusi!Benford}
Contoh distribusi berikutnya berasal dari kajian tentang digit pertama dalam himpunan data.
Ternyata banyak himpunan data yang muncul ``di dunia nyata" mempunyai sifat bahwa digit
pertama datanya tidak terdistribusi secara seragam pada himpunan $\{1, 2, \ldots, 9\}$.
Sebaliknya, digit 1 tampaknya paling mungkin muncul, dan distribusinya menurun secara
monoton pada himpunan digit yang mungkin. Dalam banyak kasus, distribusi Benford tampaknya
cocok dengan data semacam itu. Berbagai penjelasan telah diberikan mengenai munculnya
distribusi ini. Penjelasan yang mungkin paling meyakinkan adalah bahwa distribusi ini
merupakan satu-satunya distribusi yang invarian terhadap perubahan skala. Jika kita
menganggap himpunan data tertentu muncul secara ``alami'', distribusinya seharusnya tidak
dipengaruhi oleh satuan yang dipilih untuk menyatakan data tersebut; dengan kata lain,
distribusinya seharusnya invarian terhadap perubahan skala.
\par
Theodore Hill\index{HILL, T.}\footnote{T. P. Hill, ``The Significant Digit Phenomenon," \emx
{American Mathematical Monthly,} vol.\ 102, no.\ 4 (April 1995), pgs. 322-327.} memberikan
deskripsi umum distribusi Benford ketika kita meninjau $d$ digit pertama dari
bilangan-bilangan bulat dalam suatu himpunan data. Kita akan membatasi perhatian pada
digit pertama. Dalam hal ini, distribusi Benford mempunyai fungsi distribusi
$$f(k) = \log_{10}(k+1) - \log_{10}(k)\ ,$$
untuk $1 \le k \le 9$.
\par
Mark Nigrini\index{NIGRINI, M.}\footnote{M. Nigrini, ``Detecting Biases and Irregularities in
Tabulated Data," working paper} menganjurkan penggunaan distribusi Benford sebagai cara
untuk menguji catatan keuangan yang mencurigakan\index{catatan keuangan!mencurigakan},
misalnya entri pembukuan, cek, dan laporan pajak\index{laporan pajak}. Gagasannya adalah
bahwa, jika seseorang ``mengarang" bilangan-bilangan dalam kasus ini, orang tersebut
mungkin menghasilkan bilangan yang terdistribusi cukup seragam. Sebaliknya, jika bilangan
aktual digunakan, digit pertamanya akan kurang lebih mengikuti distribusi Benford. Sebagai
contoh, Nigrini menganalisis laporan pajak Presiden Clinton\index{Clinton, Bill} selama
13 tahun. Pada Gambar~\ref{fig 5.1.6}, nilai distribusi Benford ditampilkan sebagai
persegi, sedangkan data laporan pajak Presiden ditampilkan sebagai lingkaran. Dalam contoh
ini, terlihat bahwa distribusi Benford sangat cocok dengan data tersebut.
\putfig{4.5truein}{PSfig5-1-6}{Digit pertama dalam laporan pajak Presiden Clinton.}{fig 5.1.6} 
\par
Distribusi ini ditemukan oleh astronom Simon Newcomb\index{NEWCOMB, S.}, yang menulis
pernyataan berikut dalam makalahnya tentang topik tersebut: ``Bahwa kesepuluh digit tidak
muncul dengan frekuensi yang sama pasti jelas bagi siapa pun yang menggunakan tabel
logaritma dan memperhatikan bahwa halaman-halaman awal jauh lebih cepat usang daripada
halaman-halaman akhir. Angka signifikan pertama lebih sering berupa 1 daripada digit lain,
dan frekuensinya menurun hingga 9."\footnote{S. Newcomb, ``Note on the frequency of use of the different
digits in natural numbers," \emx {American Journal of Mathematics,} vol.\ 4 (1881), pgs.\ 39-40.}
\exercises
\begin{LJSItem}

\i\label{exer 5.1.100} Untuk peubah acak manakah di bawah ini distribusi seragam
layak ditetapkan?
\begin{enumerate}
\item Misalkan $X$ menyatakan hasil pelemparan satu dadu.
\item Misalkan $X$ menyatakan banyaknya sisi gambar yang diperoleh dalam tiga lemparan koin.
\item Sebuah roda rolet mempunyai 38 hasil yang mungkin: 0, 00, dan 1 sampai 36.
Misalkan $X$ menyatakan hasil ketika roda rolet diputar.
\item Misalkan $X$ menyatakan tanggal lahir seseorang yang dipilih secara acak.
\item Misalkan $X$ menyatakan banyaknya lemparan koin yang diperlukan untuk memperoleh
sisi gambar untuk pertama kalinya. 
\end{enumerate}

\i\label{exer 5.1.101} Misalkan $n$ bilangan bulat positif. Misalkan $S$ adalah
himpunan bilangan bulat dari 1 sampai $n$. Tinjau proses berikut: kita mengambil
sebuah bilangan dari $S$ secara acak dan menuliskannya. Proses ini diulangi sampai
$S$ kosong. Hasilnya adalah suatu permutasi bilangan bulat dari 1 sampai $n$.
Misalkan $X$ menyatakan permutasi ini. Apakah $X$ terdistribusi seragam?

\i\label{exer 5.1.102} Misalkan $X$ adalah peubah acak yang dapat mengambil
terhitung banyak nilai. Tunjukkan bahwa $X$ tidak mungkin terdistribusi seragam.

\i\label{exer 5.1.103} Misalkan kita berkuliah di perguruan tinggi yang mempunyai
3000 mahasiswa. Kita ingin memilih subhimpunan berukuran 100 dari seluruh mahasiswa.
Misalkan $X$ menyatakan subhimpunan yang dipilih dengan salah satu strategi berikut.
Untuk strategi manakah distribusi seragam layak ditetapkan pada $X$? Jika layak,
berapa peluang yang harus ditetapkan pada setiap hasil?
\begin{enumerate}
\item Ambil 100 mahasiswa pertama yang memasuki kafetaria untuk makan siang.
\item Minta Biro Registrasi mengurutkan mahasiswa berdasarkan nomor Jaminan Sosial,
lalu ambil 100 mahasiswa pertama pada daftar yang dihasilkan.
\item Minta Biro Registrasi menyediakan sekumpulan kartu, dengan setiap kartu memuat
nama tepat satu mahasiswa dan setiap mahasiswa muncul pada tepat satu kartu. Lemparkan
kartu-kartu itu dari jendela lantai tiga, lalu berjalanlah ke luar dan ambil 100 kartu
pertama yang ditemukan.
\end{enumerate}

\i\label{exer 5.1.104} Dengan kondisi yang sama seperti pada latihan sebelumnya,
dapatkah Anda menjelaskan suatu prosedur yang menghasilkan setiap kemungkinan hasil
dengan peluang yang sama? Dapatkah Anda menjelaskan prosedur semacam itu yang tidak
bergantung pada komputer atau kalkulator?

\i\label{exer 5.1.105} Misalkan $X_1,\ X_2,\ \ldots,\ X_n$ adalah $n$ peubah acak
yang saling bebas dan masing-masing terdistribusi seragam pada bilangan bulat dari 1
sampai $k$. Misalkan $Y$ menyatakan nilai minimum di antara semua $X_i$. Tentukan
distribusi $Y$.

\i\label{exer 5.1.16} Sebuah dadu dilempar sampai angka enam muncul untuk pertama
kalinya pada waktu $T$.
\begin{enumerate}

\item Apa distribusi peluang $T$?

\item Tentukan $P(T > 3)$.

\item Tentukan $P(T > 6 | T > 3)$.
\end{enumerate}

\i\label{5.1.18} Jika sebuah koin dilempar berulang kali, berapakah peluang bahwa
sisi gambar pertama muncul setelah lemparan kelima, dengan diketahui bahwa sisi itu
belum muncul dalam dua lemparan pertama?

\i\label{exer 5.1.106} Seorang petugas Dinas Perikanan dan Satwa Liar ditugaskan
memperkirakan jumlah ikan trout\index{ikan trout} di sebuah danau berukuran sedang.
Ia melakukan langkah berikut: ia menangkap 100 ikan trout, memberi tanda pada
masing-masing ikan, lalu mengembalikannya ke danau. Satu bulan kemudian, ia menangkap
100 ikan trout lagi dan mendapati bahwa 10 di antaranya bertanda.  
\begin{enumerate}
\item Tanpa melakukan perhitungan yang rumit, berikan perkiraan kasar jumlah ikan
trout di danau tersebut.
\item Misalkan $N$ adalah jumlah ikan trout di danau. Tentukan bentuk dalam $N$ untuk
peluang bahwa petugas tersebut menangkap 10 ikan trout bertanda di antara 100 ikan
trout yang ditangkapnya pada kesempatan kedua.
\item Tentukan nilai $N$ yang memaksimumkan bentuk pada bagian (b). Nilai ini disebut
\emx {estimasi kemungkinan maksimum}\index{estimasi kemungkinan maksimum} untuk besaran
yang tidak diketahui, yaitu $N$. \emx {Petunjuk}: Tinjau rasio bentuk untuk nilai-nilai
$N$ yang berurutan.
\end{enumerate}

\i\label{exercise 5.1.107} Sensus di Amerika Serikat merupakan upaya menghitung
setiap orang di negara tersebut. Tidak terelakkan bahwa banyak orang tidak terhitung.
Biro Sensus Amerika Serikat mengusulkan cara untuk memperkirakan jumlah orang yang
tidak terhitung dalam sensus terakhir. Usulannya sebagai berikut: di suatu wilayah,
misalkan $N$ menyatakan jumlah sebenarnya orang yang tinggal di sana. Anggap bahwa
sensus menghitung $n_1$ orang yang tinggal di wilayah ini. Kemudian sensus lain
dilakukan di wilayah tersebut dan menghitung $n_2$ orang. Selain itu, $n_{12}$ orang
terhitung pada kedua sensus.  
\begin{enumerate}
\item Dengan diketahui $N$, $n_1$, dan $n_2$, misalkan $X$ menyatakan jumlah orang
yang terhitung pada kedua sensus. Tentukan peluang bahwa $X = k$, dengan $k$ bilangan
bulat positif tetap antara 0 dan $n_2$.
\item Sekarang anggap bahwa $X = n_{12}$. Tentukan nilai $N$ yang memaksimumkan bentuk
pada bagian (a). \emx {Petunjuk}: Tinjau rasio bentuk untuk nilai-nilai $N$ yang
berurutan.
\end{enumerate}

\i\label{exer 5.1.26} Misalkan $X$ adalah peubah acak yang menyatakan banyaknya
panggilan masuk ke kantor polisi dalam selang satu menit. Dalam teks, kita telah
menunjukkan bahwa $X$ dapat dimodelkan dengan distribusi Poisson berparameter
$\lambda$, dengan parameter ini menyatakan rata-rata banyaknya panggilan masuk per
menit. Sekarang misalkan $Y$ adalah peubah acak yang menyatakan banyaknya panggilan
masuk dalam selang sepanjang $t$. Tunjukkan bahwa distribusi $Y$ diberikan oleh
$$P(Y = k) = e^{-\lambda t}{{(\lambda t)^k}\over{k!}}\ ,$$  yakni $Y$ berdistribusi
Poisson dengan parameter $\lambda t$. \emx {Petunjuk}: Misalkan seorang pengamat dari Mars
mengamati kantor polisi tersebut. Anggap pula bahwa selang waktu dasar yang digunakan
di Mars tepat sama dengan $t$ menit Bumi. Terakhir, kita akan menganggap bahwa
pengamat itu memahami penurunan distribusi Poisson dalam teks. Apa yang akan ia
tuliskan sebagai distribusi $Y$?

\i\label{exer 5.1.27} Tunjukkan bahwa nilai-nilai distribusi Poisson yang diberikan
dalam Persamaan~\ref{eq 5.1} berjumlah 1.

\i\label{exer 5.1.108} Distribusi Poisson dengan parameter $\lambda = .3$ telah
ditetapkan pada hasil suatu percobaan. Misalkan $X$ adalah fungsi hasil. Tentukan
$P(X = 0)$, $P(X = 1)$, dan $P(X > 1)$.

\i\label{exer 5.1.109} Secara rata-rata, hanya 1 dari 1000 orang yang mempunyai
golongan darah langka tertentu.

\begin{enumerate}
\item Tentukan peluang bahwa, di sebuah kota berpenduduk 10{,}000 orang, tidak ada
seorang pun yang mempunyai golongan darah ini.

\item Berapa orang harus diuji agar peluang menemukan sekurang-kurangnya satu orang
dengan golongan darah ini lebih besar daripada~1/2?
\end{enumerate}

\i\label{exer 5.1.110} Tulislah program yang memungkinkan pengguna memasukkan
$n$,~$p$,~$j$, lalu mencetak nilai eksak $b(n, p, k)$ dan aproksimasi Poisson untuk
nilai tersebut.

\i\label{exer 5.1.111} Anggap bahwa, selama setiap detik, sentral telepon Dartmouth
menerima satu panggilan dengan peluang~.01 dan tidak menerima panggilan dengan
peluang~.99. Gunakan aproksimasi Poisson untuk memperkirakan peluang bahwa operator
melewatkan paling banyak satu panggilan jika ia beristirahat minum kopi selama 5 menit.

\i\label{exer 5.1.112} Peluang memperoleh royal flush dalam satu tangan poker adalah
$p = 1/649{,}740$. Seberapa besar $n$ agar peluang tidak memperoleh royal flush dalam
$n$ tangan lebih kecil daripada $1/e$?

\i\label{exer 5.1.113} Seorang pembuat roti mencampurkan 600 kismis dan 400 keping
cokelat ke dalam adonan, lalu membuat 500 kue kering dari adonan tersebut.
\begin{enumerate}
\item Tentukan peluang bahwa kue kering yang dipilih secara acak tidak mengandung
kismis.

\item Tentukan peluang bahwa kue kering yang dipilih secara acak mengandung tepat dua
keping cokelat.

\item Tentukan peluang bahwa kue kering yang dipilih secara acak mengandung
sekurang-kurangnya dua butir (kismis atau keping cokelat).
\end{enumerate}

\i\label{exer 5.1.114} Peluang bahwa, dalam satu pembagian kartu bridge\index{bridge},
salah satu dari empat tangan seluruhnya terdiri atas kartu hati kira-kira
$6.3 \times 10^{-12}$. Di sebuah kota dengan sekitar 50{,}000 pemain bridge, pakar
peluang setempat rata-rata ditelepon sekali setahun (biasanya larut malam) dan diberi
tahu bahwa penelepon baru saja memperoleh satu tangan yang seluruhnya kartu hati.
Haruskah ia mencurigai bahwa sebagian penelepon ini menjadi korban lelucon?

\i\label{exer 5.1.115} Seorang pengiklan menjatuhkan 10{,}000 selebaran di sebuah
kota yang mempunyai 2000 blok. Anggap bahwa setiap selebaran mempunyai peluang yang
sama untuk mendarat pada setiap blok. Berapakah peluang bahwa suatu blok tertentu
tidak menerima selebaran?

\i\label{exer 5.1.116} Dalam kelas yang terdiri atas 80 mahasiswa, dosen memanggil
1~mahasiswa yang dipilih secara acak untuk menjawab pada setiap pertemuan. Terdapat
32 pertemuan dalam satu semester.
\begin{enumerate}
\item Tulislah rumus peluang eksak bahwa seorang mahasiswa tertentu dipanggil
$j$ kali selama semester tersebut.

\item Tulislah rumus aproksimasi Poisson untuk peluang ini. Dengan memakai rumus
tersebut, perkirakan peluang bahwa seorang mahasiswa tertentu dipanggil lebih dari dua kali.
\end{enumerate}

\i\label{exer 5.1.29} Anggap kita membuat kue kering kismis. Kita memasukkan sekotak
600 kismis ke dalam adonan, mencampurnya, lalu membuat 500 kue kering. Kita ingin
mengetahui peluang bahwa kue kering yang dipilih secara acak mengandung
0,~1, 2,~\dots\ kismis. Pandang kue-kue tersebut sebagai percobaan, dan misalkan $X$
adalah peubah acak yang memberikan jumlah kismis dalam suatu kue. Kita dapat
memandang jumlah kismis dalam sebuah kue sebagai hasil dari $n = 600$ percobaan saling
bebas dengan peluang keberhasilan $p = 1/500$ pada setiap percobaan. Karena $n$ besar
dan $p$ kecil, kita dapat memakai aproksimasi Poisson dengan
$\lambda = 600(1/500) = 1.2$. Tentukan peluang bahwa suatu kue mengandung
sekurang-kurangnya lima kismis.  

\i\label{exer 5.1.117} Pada suatu percobaan, distribusi Poisson dengan parameter
$\lambda = m$ telah ditetapkan. Tunjukkan bahwa salah satu hasil yang paling mungkin
adalah nilai bilangan bulat~$k$ sedemikian sehingga $m - 1 \leq k \leq m$. Dalam
kondisi apa terdapat dua nilai yang paling mungkin? \emx {Petunjuk}: Tinjau rasio
peluang-peluang yang berurutan.

\i\label{exer 5.1.118} Ketika John Kemeny\index{KEMENY, J. G.} menjabat ketua Jurusan
Matematika di Dartmouth College, ia menerima rata-rata sepuluh surat setiap hari.
Pada suatu hari kerja ia tidak menerima surat dan bertanya-tanya apakah hari itu hari
libur. Untuk memutuskannya, ia menghitung peluang bahwa dalam sepuluh tahun terdapat
sekurang-kurangnya 1~hari tanpa surat. Ia menganggap banyaknya surat yang diterima
pada suatu hari berdistribusi Poisson. Peluang berapakah yang ia peroleh?
\emx {Petunjuk}: Terapkan distribusi Poisson dua kali. Pertama, tentukan peluang bahwa
dalam 3000 hari terdapat sekurang-kurangnya 1~hari tanpa surat, dengan menganggap
setiap tahun mempunyai sekitar 300 hari pengantaran surat.
 
\i\label{exer 5.1.119} Reese Prosser\index{PROSSER, R.} tidak pernah memasukkan uang
ke meteran parkir 10 sen di Hanover. Ia menganggap peluang tertangkap adalah~.05.
Pelanggaran pertama tidak dikenai biaya, pelanggaran kedua dikenai 2~dolar, dan setiap
pelanggaran berikutnya dikenai 5~dolar. Berdasarkan anggapannya, bagaimana perbandingan
biaya harapan untuk parkir 100 kali tanpa membayar meteran dengan biaya membayar
meteran setiap kali?

\i\label{exer 9.2.15} Feller\index{FELLER, W.}\footnote{ibid., p.~161.} membahas
statistika hantaman bom terbang\index{bom terbang} di suatu wilayah di London selatan
selama Perang Dunia Kedua. Wilayah tersebut dibagi menjadi
$24 \times 24 = 576$ daerah kecil. Jumlah seluruh hantaman adalah 537. Terdapat 229
petak dengan 0~hantaman, 211 dengan 1~hantaman, 93 dengan 2~hantaman, 35 dengan
3~hantaman, 7 dengan 4~hantaman, dan 1 dengan~5 hantaman atau lebih. Dengan
menganggap hantaman-hantaman itu sepenuhnya acak, gunakan aproksimasi Poisson untuk
menentukan peluang bahwa suatu petak tertentu menerima tepat $k$~hantaman. Hitung
jumlah harapan petak yang menerima 0,~1, 2, 3, 4, dan~5 hantaman atau lebih, lalu
bandingkan dengan hasil teramati.

\i\label{exer 5.1.120} Anggap peluang terjadinya kecelakaan besar di sebuah pembangkit
listrik tenaga nuklir dalam satu tahun adalah~.001. Jika suatu negara mempunyai 100
pembangkit nuklir, perkirakan peluang terjadinya sekurang-kurangnya satu kecelakaan
semacam itu dalam suatu tahun.

\i\label{exer 5.1.121} Sebuah maskapai mendapati bahwa 4~persen penumpang yang membuat
reservasi pada penerbangan tertentu tidak akan datang. Oleh karena itu, kebijakannya
adalah menjual 100 reservasi kursi untuk pesawat yang hanya mempunyai 98~kursi.
Tentukan peluang bahwa setiap orang yang datang untuk penerbangan tersebut memperoleh
kursi.

\i\label{exer 5.1.122} Pengelola uang logam raja mengemas uang logam, 500 keping per
kotak, dan menaruh 1 uang logam palsu dalam setiap kotak. Raja curiga, tetapi alih-alih
menguji semua uang logam dalam 1~kotak, ia menguji 1~keping yang dipilih secara acak
dari masing-masing 500~kotak. Berapakah peluang ia menemukan sekurang-kurangnya satu
uang palsu? Berapakah peluangnya jika raja menguji 2~keping dari masing-masing
250 kotak?

\i\label{exer 5.1.123} (Dari Kemeny\footnote{Komunikasi pribadi.}) Tunjukkan bahwa,
jika Anda memasang 100 taruhan pada angka~17 dalam permainan rolet di Monte Carlo
(lihat Contoh~\ref{exam 6.7}), peluang Anda memperoleh keuntungan lebih besar
daripada~1/2. Berapa keuntungan harapan Anda?

\i\label{exer 9.2.20} Dalam salah satu kajian awal tentang distribusi Poisson, von
Bortkiewicz\index{von BORTKIEWICZ, L.}\footnote{L. von Bortkiewicz,  \emx {Das Gesetz der Kleinen
Zahlen} (Leipzig: Teubner, 1898), p.\ 24.} meninjau frekuensi kematian akibat
tendangan bagal\index{tendangan bagal} dalam korps tentara Prusia. Dari kajian atas
14 korps selama 20 tahun, ia memperoleh data pada Tabel~\ref{table 5.5}.
\begin{table}
\centering
\begin{tabular}{|c|c|}
\hline Jumlah kematian & Jumlah korps dengan $x$ kematian pada suatu tahun \\ \hline 
0 & 144 \\
1 & \hspace{.1in}91 \\ 
2 & \hspace{.1in}32 \\ 
3 & \hspace{.12in}11 \\ 
4 & \hspace{.18in}2 \\ \hline
\end{tabular}
\caption{Tendangan bagal.}
\label{table 5.5}
\end{table}
Sesuaikan distribusi Poisson pada data ini dan tentukan apakah menurut Anda distribusi
Poisson tersebut sesuai.

\i\label{exer 9.2.21} Sering diasumsikan bahwa lalu lintas kendaraan yang tiba di
suatu persimpangan dalam satu satuan waktu berdistribusi Poisson dengan nilai harapan
$m$. Anggap banyaknya mobil $X$ yang tiba di persimpangan dari utara dalam satuan
waktu berdistribusi Poisson dengan parameter $\lambda = m$, sedangkan banyaknya mobil
$Y$ yang tiba dari barat dalam satuan waktu berdistribusi Poisson dengan parameter
$\lambda = \bar m$. Jika $X$~dan~$Y$ saling bebas, tunjukkan bahwa jumlah seluruh
mobil $X + Y$ yang tiba di persimpangan dalam satuan waktu berdistribusi Poisson
dengan parameter $\lambda =  m + \bar m$.

\i\label{exer 9.2.22} Mobil yang melaju di Magnolia Street tiba di percabangan jalan
dan harus memilih Willow Street atau Main Street untuk melanjutkan perjalanan.
Anggap banyaknya mobil yang tiba di percabangan dalam satuan waktu berdistribusi
Poisson dengan parameter $\lambda = 4$. Sebuah mobil yang tiba di percabangan memilih
Main Street dengan peluang~3/4 dan Willow Street dengan peluang~1/4. Misalkan $X$
adalah peubah acak yang menghitung banyaknya mobil yang, dalam suatu satuan waktu,
melewati Joe's Barber Shop di Main Street. Apa distribusi $X$?

\i\label{exer 9.2.23} Dalam banding perkara \emx {People v.\ Collins}\index{Collins, People
v.}\index{People v. Collins} (lihat Latihan~\ref{sec 4.1}.\ref{exer 4.1.26}),
kuasa hukum pembela berargumen sebagai berikut: misalkan, sebagai contoh, terdapat
5{,}000{,}000 pasangan di wilayah Los Angeles dan peluang bahwa pasangan yang dipilih
secara acak sesuai dengan deskripsi para saksi adalah 1/12{,}000{,}000. Maka peluang
terdapat dua pasangan semacam itu, dengan diketahui terdapat sekurang-kurangnya satu,
sama sekali tidak kecil. Tentukan peluang ini. (Mahkamah Agung California membatalkan
putusan bersalah semula.)


\i\label{exer 5.1.20} Suatu lot produksi baut penegang kuningan berisi $S$ barang,
dengan $D$ di antaranya cacat. Sampel sebanyak $s$ barang diambil tanpa pengembalian.
Misalkan $X$ adalah peubah acak yang memberikan banyaknya barang cacat dalam sampel.
Misalkan $p(d) = P(X = d)$.

\begin{enumerate}
\item Tunjukkan bahwa
$$ p(d) = \frac{{D \choose d} {{S - D} \choose {s - d}}}{{S \choose s}}\ .
$$ Dengan demikian, X berdistribusi hipergeometrik.

\item Buktikan identitas berikut, yang dikenal sebagai \emx {rumus Euler}\index{rumus Euler}:
$$
\sum_{d = 0}^{\min(D,s)}{ D \choose d}   {{S - D} \choose {s - d}} =  {S \choose s}\ .
$$
\end{enumerate}

\i\label{exer 5.1.21} Sebuah wadah berisi 1000 baut penegang dengan jumlah barang
cacat $D$ yang tidak diketahui. Dalam sampel 100 baut penegang terdapat 2 barang
cacat. \emx {Estimasi kemungkinan maksimum}\index{estimasi kemungkinan maksimum}
untuk $D$ adalah jumlah barang cacat yang memberikan peluang tertinggi untuk memperoleh
jumlah barang cacat yang teramati dalam sampel. Tebaklah nilai $D$, lalu tulislah
program komputer untuk memverifikasi tebakan Anda.

\i\label{exer 5.1.22} Terdapat rusa besar\index{rusa besar} dalam jumlah yang tidak
diketahui di Isle Royale\index{Isle Royale} (sebuah taman nasional di Danau Superior).
Untuk memperkirakan jumlah rusa besar, 50 ekor ditangkap dan diberi tanda. Enam bulan
kemudian, 200 ekor ditangkap dan ternyata 8 di antaranya bertanda. Perkirakan jumlah
rusa besar di Isle Royale dari data ini, lalu verifikasi tebakan Anda dengan program
komputer (lihat Latihan~\ref{exer 5.1.21}).

\i\label{exer 5.1.23} Suatu lot produksi cambuk kereta kuda berisi 20 barang, dengan
5 di antaranya cacat. Sampel acak sebanyak~5 barang dipilih untuk diperiksa. Tentukan
peluang bahwa sampel tersebut memuat tepat satu barang cacat
\begin{enumerate}
\item jika pengambilan sampel dilakukan dengan pengembalian.

\item jika pengambilan sampel dilakukan tanpa pengembalian.
\end{enumerate}

\i\label{exer 5.1.28} Misalkan $N$ dan $k$ menuju $\infty$ sedemikian rupa sehingga
$k/N$ tetap. Tunjukkan bahwa
$$h(N, k, n, x) \rightarrow b(n, k/N, x)\ .$$

\i\label{exer 5.1.24} Satu dek bridge\index{bridge} berisi 52 kartu, dengan 13 kartu
dalam masing-masing dari empat jenis: sekop, hati, wajik, dan keriting. Satu tangan
berisi 13 kartu dibagikan dari dek yang telah dikocok. Tentukan peluang bahwa tangan
tersebut mempunyai
\begin{enumerate}

\item distribusi jenis kartu 4, 4, 3, 2 (misalnya empat sekop, empat hati, tiga wajik,
dan dua keriting).

\item distribusi jenis kartu 5, 3, 3, 2.
\end{enumerate}

\i\label{exer 5.1.125} Tulislah algoritma komputer yang menyimulasikan peubah acak
hipergeometrik dengan parameter $N$,  $k$, dan $n$.

\i\label{exer 7.1.9} Anda diberi empat dadu yang berbeda. Dadu pertama mempunyai dua
sisi bertanda~0 dan empat sisi bertanda~4. Dadu kedua bertanda 3 pada setiap sisinya.
Dadu ketiga bertanda 2 pada empat sisi dan 6 pada dua sisi, sedangkan dadu keempat
bertanda 1 pada tiga sisi dan 5 pada tiga sisi. Anda mempersilakan teman memilih salah
satu dari keempat dadu tersebut. Kemudian Anda memilih salah satu dari tiga dadu yang
tersisa, dan masing-masing melempar dadu pilihan. Orang yang memperoleh angka terbesar
memenangkan satu dolar. Tunjukkan bahwa Anda dapat memilih dadu sehingga peluang Anda
menang adalah 2/3, apa pun dadu yang dipilih teman. (Lihat Tenney dan
Foster.\footnote{R.~L.~Tenney and C.~C.~Foster,  \emx 
{Non-transitive Dominance}, Math. Mag. 49 (1976) no. 3, pgs. 115-120.})

\i\label{exer 5.1.25} Mahasiswa dalam suatu kelas dikelompokkan berdasarkan warna
rambut dan warna mata. Konvensi yang digunakan adalah: rambut cokelat dan hitam
dianggap gelap, sedangkan rambut merah dan pirang dianggap terang; mata hitam dan
cokelat dianggap gelap, sedangkan mata biru dan hijau dianggap terang. Mereka
mengumpulkan data yang ditampilkan pada Tabel~\ref{table 5.6}.  
\begin{table}[t]
\centering
\begin{tabular}{|l|c|c|c|}
\hline
 & Mata Gelap & Mata Terang & \\ \hline 
Rambut Gelap  & 28              & 15 & \hspace{.25in} 43 \hspace{.25in}   \\ \hline 
Rambut Terang & \hspace{.1in}9  & 23 & \hspace{.25in} 32 \hspace{.25in}   \\ \hline
           & 37              & 38 & \hspace{.25in} 75 \hspace{.25in}   \\ \hline
\end{tabular}
\caption{Data teramati.}
\label{table 5.6}
\end{table}
Apakah kedua ciri ini saling bebas? (Lihat Contoh~\ref{exam 5.6}.)

\i\label{exer 5.1.124} Misalkan dalam distribusi hipergeometrik, $N$ dan $k$ menuju
$\infty$ sedemikian rupa sehingga rasio $k/N$ mendekati bilangan real $p$ antara
0 dan 1. Tunjukkan bahwa distribusi hipergeometrik menuju distribusi binomial dengan
parameter $n$ dan $p$.  

\i\label{exer 5.1.126}  
\begin{enumerate}
\item Hitung digit pertama dari 100 pangkat pertama bilangan 2, lalu tentukan seberapa
baik data ini sesuai dengan distribusi Benford.
\item Kalikan setiap bilangan dalam himpunan data pada bagian (a) dengan 3, lalu
bandingkan distribusi digit pertamanya dengan distribusi Benford.
\end{enumerate}

\i\label{exer 5.1.127} Dalam lotre Powerball\index{lotre Powerball}\index{lotre!Powerball},
peserta memilih 5 bilangan bulat berbeda antara 1 dan 45, lalu memilih satu bilangan bonus dari
rentang yang sama (bilangan bonus boleh sama dengan salah satu dari lima bilangan pertama).
Sebagian peserta memilih sendiri bilangannya, sedangkan yang lain membiarkan komputer
memilih. Data pada Tabel~\ref{table 5.10} adalah bilangan-bilangan yang dipilih peserta
di suatu negara bagian pada 3 Mei 1996. Grafik pancang data tersebut ditampilkan pada
Gambar~\ref{fig 5.2.1}.  
\putfig{4truein}{PSfig5-2-1}{Distribusi pilihan dalam lotre Powerball.}{fig 5.2.1}
\par
\choice{Menurut Anda, apakah orang-orang memilih bilangan secara acak? Jelaskan jawaban Anda.}{Tujuan
soal ini adalah menguji hipotesis bahwa bilangan yang dipilih terdistribusi seragam. Untuk itu,
hitung nilai $v$ dari peubah acak $\chi^2$ yang diberikan dalam
Contoh~\ref{exam 5.6}. Dalam kasus ini, peubah acak tersebut mempunyai 44 derajat
kebebasan. Dari tabel $\chi^2$, kita dapat menemukan nilai $v_0 = 59.43$, yaitu bilangan
dengan sifat bahwa peubah acak berdistribusi $\chi^2$ mengambil nilai yang melebihi $v_0$
hanya 5\% dari waktu. Apakah nilai $v$ yang Anda hitung melebihi $v_0$? Jika demikian,
Anda harus menolak hipotesis bahwa pilihan peserta terdistribusi seragam.}

\begin{table}[h]
\centering
\begin{tabular}{llllll} Bilangan & Banyaknya & Bilangan & Banyaknya & Bilangan & Banyaknya\\
&Dipilih && Dipilih &&Dipilih\\
\hline
1&2646&
2&2934&
3&3352\\
4&3000&
5&3357&
6&2892\\
7&3657&
8&3025&
9&3362\\
10&2985&
11&3138&
12&3043\\
13&2690&
14&2423&
15&2556\\
16&2456&
17&2479&
18&2276\\
19&2304&
20&1971&
21&2543\\
22&2678&
23&2729&
24&2414\\
25&2616&
26&2426&
27&2381\\
28&2059&
29&2039&
30&2298\\
31&2081&
32&1508&
33&1887\\
34&1463&
35&1594&
36&1354\\
37&1049&
38&1165&
39&1248\\
40&1493&
41&1322&
42&1423\\
43&1207&
44&1259&
45&1224\\
\end{tabular}
\caption{Bilangan yang dipilih peserta lotre Powerball.}
\label{table 5.10}
\end{table}


\end{LJSItem}

%Again, the choice function seems to be screwing up.
%\choice{}{\section{Important Densities}\label{sec 5.2}
\section{Densitas Penting}\label{sec 5.2}
Dalam bagian ini, kita akan memperkenalkan beberapa fungsi densitas peluang yang
penting dan memberikan beberapa contoh penggunaannya.  Kita juga akan membahas
cara menyimulasikan suatu densitas dengan komputer.

\subsection*{Densitas Seragam Kontinu}\index{fungsi densitas!seragam}\index{densitas seragam}

Fungsi densitas yang paling sederhana berkaitan dengan peubah acak $U$ yang nilainya
menyatakan hasil percobaan memilih sebuah bilangan real secara acak dari interval
$[a, b]$.  
$$ f(\omega) = \left \{ \matrix{ 
                       1/(b - a), &\,\,\, \mbox{jika}\,\,\, a \leq \omega \leq b, \cr
                       0,         &\,\,\, \mbox{selainnya.}\cr}\right. 
$$

Kita dapat menyimulasikan densitas ini dengan mudah pada komputer.  Kita cukup
menghitung ekspresi
$$ (b - a) rnd + a\ .
$$

\subsection*{Densitas Eksponensial dan Gamma}\index{densitas eksponensial}\index{fungsi
densitas!eksponensial}  Fungsi densitas eksponensial didefinisikan oleh
$$ f(x) = \left \{ \matrix{
                   \lambda e^{-\lambda x}, &\,\,\, \mbox{jika}\,\,\, 0 \leq x < \infty, \cr
                   0,                      &\,\,\, \mbox{selainnya}. \cr} \right.
$$ Di sini $\lambda$ adalah sebarang konstanta positif yang bergantung pada
percobaannya.  Pembaca telah menjumpai densitas ini dalam Contoh~\ref{exam 2.2.7.5}.
Dalam Gambar~\ref{fig 2.20}, kita menampilkan grafik beberapa densitas eksponensial
untuk berbagai pilihan
$\lambda$.  
\putfig{4.5truein}{PSfig2-20}{Densitas eksponensial.}{fig 2.20} 
Densitas eksponensial sering digunakan untuk menjelaskan percobaan yang melibatkan
pertanyaan: Berapa lama sampai sesuatu terjadi?  Sebagai contoh, densitas eksponensial
sering digunakan untuk mempelajari waktu di antara emisi partikel dari sumber
radioaktif.
\par Fungsi distribusi kumulatif densitas eksponensial mudah dihitung.  Misalkan $T$
adalah peubah acak berdistribusi eksponensial dengan parameter $\lambda$.  Jika
$x \ge 0$, maka
\begin{eqnarray*} F(x) & = & P(T \le x) \\ & = & \int_0^x \lambda e^{-\lambda t}\,dt
\\ & = & 1 - e^{-\lambda x}\ .\\
\end{eqnarray*}
\par
Baik densitas eksponensial maupun distribusi geometrik memiliki sifat yang dikenal
sebagai sifat ``tanpa memori"\index{sifat tanpa memori}.  Sifat ini diperkenalkan
dalam Contoh~\ref{exam 5.1}; sifat tersebut menyatakan bahwa
$$P(T > r + s\,|\,T > r) = P(T > s)\ .$$ 
Kita dapat menunjukkan bahwa sifat ini berlaku bagi densitas eksponensial dengan
menghitung kedua ruas persamaan tersebut.  Ruas kanan hanyalah
$$1 - F(s) = e^{-\lambda s}\ ,$$ sedangkan ruas kiri adalah
\begin{eqnarray*} {{P(T > r + s)}\over{P(T > r)}} & = & {{1 - F(r + s)}\over{1 -
F(r)}} \\ & = & {{e^{-\lambda (r+s)}}\over{e^{-\lambda r}}} \\ & = & e^{-\lambda s}\
.\\
\end{eqnarray*}
\par
Terdapat hubungan yang sangat penting antara densitas eksponensial dan distribusi
Poisson.  Mula-mula kita definisikan $X_1,\ X_2,\ \ldots$ sebagai barisan peubah acak
independen yang masing-masing berdistribusi eksponensial dengan parameter $\lambda$.
Kita dapat memandang $X_i$ sebagai lamanya waktu antara emisi partikel ke-$i$ dan
ke-$(i+1)$ dari sebuah sumber radioaktif.  (Seperti yang akan kita lihat dalam
Bab~\ref{chp
6}, parameter $\lambda$ dapat dipandang sebagai kebalikan dari rata-rata
lamanya waktu di antara emisi.  Parameter ini merupakan besaran yang dapat diukur
dalam percobaan nyata semacam itu.)  
\par Sekarang kita tinjau interval waktu sepanjang $t$, dan kita misalkan $Y$ sebagai
peubah acak yang menghitung banyaknya emisi yang terjadi dalam interval tersebut.
Kita ingin menghitung fungsi distribusi $Y$ (jelas bahwa $Y$ adalah peubah acak
diskret).  Jika $S_n$ menyatakan jumlah
$X_1 + X_2 + 
\cdots + X_n$, mudah dilihat bahwa
$$P(Y = n) = P(S_n \le t\ \mbox{dan}\ S_{n+1} > t)\ .$$ 
Karena kejadian $S_{n+1} \le t$ merupakan himpunan bagian dari kejadian $S_n \le t$,
peluang di atas sama dengan
\begin{equation} P(S_n \le t) - P(S_{n+1} \le t)\ .\label{eq 5.8}
\end{equation} Dalam Bab~\ref{chp 7}, kita akan menunjukkan bahwa densitas $S_n$
diberikan oleh rumus berikut:
$$ g_n(x) = \left \{ \begin{array}{ll}
                       \lambda{{(\lambda x)^{n-1}}\over{(n-1)!}}e^{-\lambda x}, 
                                          & \mbox{jika $x > 0$,} \\
                               0,         & \mbox{selainnya.}
                  \end{array}
         \right. 
$$ Densitas ini merupakan contoh densitas gamma\index{densitas gamma}\index{fungsi densitas!gamma}
dengan parameter
$\lambda$ dan
$n$.  Densitas gamma umum memperbolehkan $n$ berupa sebarang bilangan real positif.
Kita tidak akan membahas densitas umum ini.
\par
Dengan induksi pada $n$, mudah ditunjukkan bahwa fungsi distribusi kumulatif $S_n$
diberikan oleh
$$ G_n(x) = \left \{ \begin{array}{ll}
                         1 - e^{-\lambda x}\biggl(1 + {{\lambda x}\over {1!}} + \cdots
+
                         {{(\lambda x)^{n-1}}\over{(n-1)!}}\biggr), & \mbox{jika $x >
0$}, \\
                         0,         & \mbox{selainnya.}
                   \end{array}
          \right. 
$$ Dengan menggunakan ekspresi ini, besaran dalam (\ref{eq 5.8}) mudah dihitung; kita
memperoleh
$$ e^{-\lambda t}{{(\lambda t)^n}\over{n!}}\ ,$$ yang akan dikenali pembaca sebagai
peluang bahwa peubah acak berdistribusi Poisson dengan parameter $\lambda t$
bernilai $n$.
\par Hubungan di atas memungkinkan kita menyimulasikan distribusi Poisson setelah
kita menemukan cara menyimulasikan densitas eksponensial.  Peubah acak berikut
melakukannya:
\begin{equation} Y = -{1\over\lambda} \log(rnd)\ .\label{eq 5.9}
\end{equation} Dengan Korolari~\ref{cor 5.2} (di bawah), kita dapat menurunkan
ekspresi di atas (lihat Latihan~\ref{exer 5.2.2.5}).  Untuk saat ini, kita cukup
memberikan perhitungan singkat yang menunjukkan bahwa peubah acak $Y$ memiliki sifat
yang diperlukan.  Kita memperoleh
\begin{eqnarray*} P(Y \le y) & = & P\Bigl(-{1\over\lambda} \log(rnd) \le y\Bigr) \\ &
= & P(\log(rnd)
\ge -\lambda y) \\ & = & P(rnd \ge e^{-\lambda y}) \\ & = & 1 - e^{-\lambda y}\ . \\
\end{eqnarray*} Ekspresi terakhir ini merupakan fungsi distribusi kumulatif peubah
acak berdistribusi eksponensial dengan parameter $\lambda$.
\par Untuk menyimulasikan peubah acak Poisson $W$ dengan parameter $\lambda$, kita
cukup menghasilkan barisan nilai peubah acak berdistribusi eksponensial dengan
parameter yang sama dan mencatat subtotal $S_k$ dari nilai-nilai tersebut.  Kita
berhenti menghasilkan barisan ketika subtotal untuk pertama kalinya melebihi
$\lambda$.  Misalkan kita mendapati bahwa
$$S_n \le \lambda < S_{n+1}\ .$$ Maka nilai $n$ dikembalikan sebagai nilai simulasi
untuk $W$.


\begin{example}(Antrean)\index{antrean}\label{exam 5.21} Misalkan pelanggan tiba pada waktu
acak di sebuah tempat pelayanan dengan satu pelayan.  Setiap pelanggan langsung dilayani jika
tidak ada orang di depannya, tetapi jika ada, ia harus menunggu gilirannya dalam antrean.  Berapa
lama waktu tunggu yang dapat diharapkan setiap pelanggan?  (Kita definisikan waktu tunggu seorang
pelanggan sebagai lamanya waktu antara kedatangan dan awal pelayanannya.)
\par
Kita asumsikan bahwa waktu antar-kedatangan pelanggan yang berturut-turut diberikan oleh peubah
acak $X_1$,~$X_2$, \dots,~$X_n$ yang saling independen dan berdistribusi identik dengan fungsi
distribusi kumulatif eksponensial
$$ F_X(t) = 1 - e^{-\lambda t}.
$$ 
Kita juga asumsikan bahwa waktu pelayanan pelanggan yang berturut-turut diberikan oleh peubah acak
$Y_1$,~$Y_2$, \dots,~$Y_n$ yang juga saling independen dan berdistribusi identik dengan fungsi
distribusi kumulatif eksponensial lain,
$$ F_Y(t) = 1 - e^{-\mu t}.
$$ 
\par
Parameter $\lambda$ dan $\mu$ masing-masing menyatakan kebalikan dari rata-rata waktu antar-
kedatangan\index{waktu antar-kedatangan, rata-rata} pelanggan dan rata-rata waktu
pelayanan\index{waktu pelayanan, rata-rata} pelanggan.  Jadi, misalnya, makin besar nilai
$\lambda$, makin kecil rata-rata waktu antar-kedatangan pelanggan.  Kita dapat menduga bahwa lama
waktu yang dihabiskan seorang pelanggan dalam antrean bergantung pada perbandingan antara rata-rata
waktu antar-kedatangan dan rata-rata waktu pelayanan.
\par
Kita dapat memeriksa dugaan ini dengan mudah melalui simulasi.  Program {\bf Queue}\index{Queue
(program)} menyimulasikan proses antrean ini.  Misalkan $N(t)$ adalah banyaknya pelanggan dalam
antrean pada waktu $t$.  Kemudian kita plot $N(t)$ sebagai fungsi $t$ untuk berbagai pilihan parameter
$\lambda$ dan
$\mu$ (lihat Gambar~\ref{fig 5.17}).

\putfig{4.9truein}{PSfig5-17}{Panjang antrean.}{fig 5.17}

Kita perhatikan bahwa ketika $\lambda < \mu$, berlaku $1/\lambda > 1/\mu$, sehingga rata-rata waktu
antar-kedatangan lebih besar daripada rata-rata waktu pelayanan; artinya, secara rata-rata pelanggan
dilayani lebih cepat daripada kedatangan pelanggan baru.  Jadi, dalam hal ini wajar untuk mengharapkan
bahwa $N(t)$ tetap kecil.  Namun, jika $\lambda > \mu$, pelanggan tiba lebih cepat daripada mereka
dilayani dan, sesuai dugaan, $N(t)$ tampak bertambah tanpa batas.
\par
Sekarang kita dapat bertanya: Berapa lama seorang pelanggan harus menunggu pelayanan dalam antrean?
Untuk menyelidiki pertanyaan ini, misalkan $W_i$ adalah lamanya waktu pelanggan ke-$i$ berada dalam
sistem (menunggu dalam antrean dan dilayani).  Dengan program {\bf Queue}, kita dapat menyajikan data
ini dalam grafik batang untuk memperoleh gambaran tentang distribusi $W_i$ (lihat
Gambar~\ref{fig 5.18}). (Di sini $\lambda = 1$ dan $\mu = 1.1$.)
\putfig{4truein}{PSfig5-18}{Waktu tunggu.}{fig 5.18}
\par
Kita melihat bahwa waktu tunggu ini tampak berdistribusi eksponensial.  Hal ini selalu berlaku
ketika $\lambda < \mu$.  Pembuktian fakta ini terlalu rumit untuk diberikan di sini, tetapi kita
dapat memeriksanya melalui simulasi untuk berbagai pilihan $\lambda$ dan $\mu$, seperti di atas.
\end{example}

\subsection*{Fungsi dari Peubah Acak}\index{peubah acak!fungsi dari} 
Sebelum melanjutkan daftar densitas penting, kita berhenti sejenak untuk meninjau peubah acak yang
merupakan fungsi dari peubah acak lain.  Kita akan membuktikan teorema umum yang memungkinkan kita
menurunkan ekspresi seperti Persamaan~\ref{eq 5.9}.
\par
\begin{theorem}\label{thm 5.1} Misalkan $X$ adalah peubah acak kontinu dan $\phi(x)$
merupakan fungsi yang naik tegas pada daerah nilai $X$.  Definisikan $Y = \phi(X)$.
Misalkan $X$ dan $Y$ masing-masing memiliki fungsi distribusi kumulatif $F_X$ dan
$F_Y$.  Kedua fungsi ini berhubungan melalui
$$ F_Y(y) = F_X(\phi^{-1}(y)).
$$ Jika $\phi(x)$ turun tegas pada daerah nilai $X$, maka
$$F_Y(y) = 1 - F_X(\phi^{-1}(y))\ .$$
\proof Karena $\phi$ naik tegas pada daerah nilai $X$, kejadian
$(X \le \phi^{-1}(y))$ dan $(\phi(X) \le y)$ sama.  Oleh karena itu,
\begin{eqnarray*} F_Y(y) & = & P(Y \le y) \\ & = & P(\phi(X) \le y) \\ & = & P(X \le
\phi^{-1}(y)) \\ & = & F_X(\phi^{-1}(y))\ . \\
\end{eqnarray*}
\par Jika $\phi(x)$ turun tegas pada daerah nilai $X$, maka
\begin{eqnarray*} F_Y(y) & = & P(Y \leq y) \\
       & = & P(\phi(X) \leq y) \\
       & = & P(X \geq \phi^{-1}(y)) \\
       & = & 1 - P(X < \phi^{-1}(y)) \\
       & = & 1 - F_X(\phi^{-1}(y))\ . \\
\end{eqnarray*} Ini menyelesaikan pembuktian.
\end{theorem}
\begin{corollary}\label{cor 5.1} Misalkan $X$ adalah peubah acak kontinu dan $\phi(x)$
merupakan fungsi yang naik tegas pada daerah nilai $X$.  Definisikan $Y = \phi(X)$.
Misalkan fungsi densitas $X$ dan $Y$ masing-masing adalah $f_X$ dan $f_Y$.  Kedua
fungsi ini berhubungan melalui
$$f_Y(y) = f_X(\phi^{-1}(y)){{d\ }\over{dy}}\phi^{-1}(y)\ .$$  Jika $\phi(x)$
turun tegas pada daerah nilai $X$, maka
$$f_Y(y) = -f_X(\phi^{-1}(y)){{d\ }\over{dy}}\phi^{-1}(y)\ .$$
\proof Hasil ini mengikuti dari Teorema~\ref{thm 5.1} dengan menggunakan Aturan
Rantai.
\end{corollary} 
\par Jika fungsi $\phi$ tidak naik tegas maupun turun tegas, situasinya agak lebih
rumit, tetapi dapat ditangani dengan metode yang sama. Sebagai contoh, misalkan
$Y = X^2$. Maka $\phi(x) = x^2$, dan
\begin{eqnarray*} F_Y(y) & = & P(Y \leq y) \\
       & = & P(-\sqrt y \leq X \leq +\sqrt y) \\
       & = & P(X \leq +\sqrt y) - P(X \leq -\sqrt y) \\
       & = & F_X(\sqrt y) - F_X(-\sqrt y)\ .\\
\end{eqnarray*} Moreover,
\begin{eqnarray*} f_Y(y) & = & \frac d{dy} F_Y(y) \\
       & = & \frac d{dy} (F_X(\sqrt y) - F_X(-\sqrt y)) \\
       & = & \Bigl(f_X(\sqrt y) + f_X(-\sqrt y)\Bigr) \frac 1{2\sqrt y}\ . \\
\end{eqnarray*}
\par 
Kita melihat bahwa untuk menyatakan $F_Y$ dalam $F_X$ ketika $Y = \phi(X)$, kita
harus menyatakan $P(Y \leq y)$ dalam $P(X \leq x)$, dan secara umum proses ini
bergantung pada struktur $\phi$.
\subsection*{Simulasi}\index{menyimulasikan peubah acak}
Teorema~\ref{thm 5.1}, antara lain, memberi tahu kita cara menyimulasikan pada
komputer sebuah peubah acak $Y$ dengan fungsi distribusi kumulatif $F$ yang telah
ditentukan.  Kita asumsikan bahwa $F(y)$ naik tegas untuk nilai-nilai $y$ yang
memenuhi $0 < F(y) < 1$.  Untuk keperluan ini, misalkan $U$ adalah peubah acak yang
berdistribusi seragam pada $[0, 1]$. Maka $U$ memiliki fungsi distribusi kumulatif
$F_U(u) = u$.  Sekarang, jika $F$ adalah fungsi distribusi kumulatif yang ditentukan
untuk $Y$, untuk menuliskan $Y$ dalam $U$ mula-mula kita selesaikan persamaan
$$ 
F(y) = u
$$ 
untuk $y$ dalam $u$.  Kita memperoleh $y = F^{-1}(u)$.  Perhatikan bahwa karena $F$
merupakan fungsi naik, persamaan ini selalu memiliki solusi tunggal (lihat
Gambar~\ref{fig
5.13}).  Selanjutnya kita tetapkan $Z = F^{-1}(U)$ dan, berdasarkan
Teorema~\ref{thm 5.1}, memperoleh
$$ 
F_Z(y) = F_U(F(y)) = F(y)\ ,
$$ 
karena $F_U(u) = u$.  Jadi, $Z$ dan $Y$ memiliki fungsi distribusi kumulatif yang
sama.  Ringkasnya, kita memperoleh hasil berikut.

\putfig{3truein}{PSfig5-13}
{Mengubah distribusi seragam $F_{U}$ menjadi distribusi $F_{Y}$ yang ditentukan.}{fig
5.13}%%4.5truein 

\begin{corollary}\label{cor 5.2}  Jika $F(y)$ adalah fungsi distribusi kumulatif yang
diberikan dan naik tegas ketika $0 < F(y) < 1$, serta jika $U$ adalah peubah acak
yang berdistribusi seragam pada $[0,1]$, maka
$$ 
Y = F^{-1}(U)
$$ 
memiliki distribusi kumulatif $F(y)$.
\end{corollary}

Jadi, untuk menyimulasikan peubah acak dengan distribusi kumulatif $F$ yang diberikan,
kita cukup menetapkan $Y = F^{-1}(\mbox{rnd})$.

\subsection*{Densitas Normal}\index{densitas normal}\index{fungsi densitas!normal} 
Sekarang kita sampai pada fungsi densitas yang paling penting, yaitu fungsi densitas
normal.  Kita telah melihat dalam Bab~\ref{chp 3} bahwa fungsi distribusi binomial
berbentuk lonceng, bahkan untuk nilai $n$ berukuran sedang.  Kita ingat bahwa peubah
acak berdistribusi binomial dengan parameter $n$ dan $p$ dapat dipandang sebagai
jumlah dari $n$ peubah acak 0-1 yang saling independen.  Sebuah teorema yang sangat
penting dalam teori peluang, yang disebut Teorema Limit Pusat, menyatakan bahwa dalam
kondisi yang sangat umum, jika kita menjumlahkan banyak peubah acak yang saling
independen, distribusi jumlahnya dapat dihampiri dengan baik oleh suatu densitas
kontinu tertentu yang disebut densitas normal.  Teorema ini akan dibahas dalam
Bab~\ref{chp 9}.
\par
Fungsi densitas normal dengan parameter $\mu$ dan $\sigma$ didefinisikan sebagai berikut:
$$
f_X(x) = \frac 1{\sqrt{2\pi}\sigma} e^{-(x - \mu)^2/2\sigma^2}\ .
$$
Parameter $\mu$ menyatakan ``pusat" densitas (dan dalam Bab~\ref{chp 6}, kita akan
menunjukkan bahwa parameter ini merupakan nilai rata-rata, atau nilai harapan,
densitas tersebut).  Parameter $\sigma$ mengukur ``sebaran" densitas, sehingga
nilainya diasumsikan positif.  (Dalam Bab~\ref{chp 6}, kita akan menunjukkan bahwa
$\sigma$ merupakan simpangan baku densitas.)  Sama sekali tidak jelas secara langsung
bahwa fungsi di atas adalah suatu densitas, yakni bahwa integralnya pada seluruh garis
real sama dengan 1.  Fungsi distribusi kumulatif diberikan oleh rumus
$$ F_X(x) = \int_{-\infty}^x \frac 1{\sqrt{2\pi}\sigma} e^{-(u -
\mu)^2/2\sigma^2}\,du\ .
$$

Dalam Gambar~\ref{fig 5.12}, sebagai perbandingan kita sertakan plot densitas normal
untuk kasus $\mu = 0$ dan $\sigma = 1$, serta $\mu = 0$ dan $\sigma = 2$.

\putfig{3truein}{PSfig5-12}{Densitas normal untuk dua himpunan nilai parameter.}{fig 5.12}%%4.5truein 

\par 
Kita tidak dapat menuliskan $F_X$ dalam fungsi-fungsi sederhana.  Hal ini menimbulkan
beberapa persoalan.  Pertama, nilai $F_X$ harus dihitung menggunakan integrasi
numerik.  Tersedia tabel-tabel lengkap yang memuat nilai fungsi ini (lihat Lampiran A).
Kedua, kita tidak dapat menuliskan $F^{-1}_X$ dalam bentuk tertutup, sehingga kita
tidak dapat menggunakan Korolari~\ref{cor 5.2} untuk membantu menyimulasikan peubah
acak normal.  Karena itu, telah dikembangkan metode-metode khusus untuk menyimulasikan
distribusi normal.  Salah satu metode tersebut bertumpu pada fakta bahwa jika $U$ dan
$V$ adalah peubah acak independen dengan densitas seragam pada $[0,1]$, maka peubah acak
$$  X = \sqrt{-2\log U} \cos 2\pi V
$$ dan
$$ Y = \sqrt{-2\log U} \sin 2\pi V
$$ saling independen dan memiliki fungsi densitas normal dengan parameter $\mu = 0$
dan $\sigma = 1$.  (Hal ini tidak langsung terlihat, dan kita juga tidak akan
membuktikannya di sini.  Lihat Box dan
Muller.\index{BOX, G. E. P.}\index{MULLER, M. E.}\footnote{G. E. P. Box and M. E. Muller,  \emx
{A Note on the Generation of Random Normal Deviates}, Ann. of Math. Stat. 29 (1958), pgs.
610-611.})
\par Misalkan $Z$ adalah peubah acak normal dengan parameter $\mu = 0$ dan
$\sigma = 1$. Peubah acak normal dengan parameter ini disebut peubah acak normal
\emx {baku}\index{peubah acak normal\\ baku}.  Fakta penting dan berguna menyatakan
bahwa jika kita menuliskan
$$X = \sigma Z + \mu\ ,$$
peubah $X$ merupakan peubah acak normal dengan parameter $\mu$ dan $\sigma$.  Untuk
menunjukkannya, kita akan menggunakan Teorema~\ref{thm 5.1}.  Kita mempunyai
$\phi(z) = \sigma z + \mu$,
$\phi^{-1}(x) = (x - \mu)/\sigma$, dan
\begin{eqnarray*} F_X(x) & = & F_Z\left(\frac {x - \mu}\sigma \right), \\ f_X(x) & = &
f_Z\left(\frac {x - \mu}\sigma \right) \cdot \frac 1\sigma \\
       & = & \frac 1{\sqrt{2\pi}\sigma} e^{-(x - \mu)^2/2\sigma^2}\ . \\
\end{eqnarray*}
Pembaca akan melihat bahwa ekspresi terakhir ini adalah fungsi densitas dengan
parameter $\mu$ dan $\sigma$, seperti yang dinyatakan.
\par
Kita telah melihat di atas bahwa peubah acak normal baku $Z$ dapat disimulasikan.
Jika kita ingin menyimulasikan peubah acak normal $X$ dengan parameter $\mu$ dan
$\sigma$, kita cukup mentransformasi nilai simulasi $Z$ dengan persamaan
$X = \sigma Z + \mu$.
\par
Misalkan kita ingin menghitung nilai fungsi distribusi kumulatif peubah acak normal
$X$ dengan parameter $\mu$ dan $\sigma$.  Perhitungan ini dapat kita reduksi menjadi
perhitungan tentang peubah acak normal baku $Z$ sebagai berikut:
\begin{eqnarray*} F_X(x) & = & P(X \leq x) \\
       & = & P\left(Z \leq \frac {x - \mu}\sigma \right) \\
       & = & F_Z\left(\frac {x - \mu}\sigma \right)\ . \\
\end{eqnarray*}
Ekspresi terakhir ini dapat ditemukan dalam tabel nilai fungsi distribusi kumulatif
peubah acak normal baku.  Jadi, kita melihat bahwa tidak perlu membuat tabel fungsi
distribusi normal untuk sebarang $\mu$ dan $\sigma$.
\par
Proses mengubah peubah acak normal menjadi peubah acak normal baku dikenal sebagai
pembakuan.  Jika $X$ berdistribusi normal dengan parameter $\mu$ dan $\sigma$, serta
jika
$$ Z = \frac{X - \mu}\sigma\ ,$$
maka $Z$ disebut versi baku dari $X$.
\par
Contoh berikut menunjukkan cara menggunakan versi baku dari peubah acak normal $X$
untuk menghitung peluang tertentu yang berkaitan dengan $X$.

\begin{example}\label{exam 5.16}
Misalkan $X$ adalah peubah acak berdistribusi normal dengan parameter $\mu = 10$ dan
$\sigma = 3$.  Carilah peluang bahwa $X$ berada di antara 4 dan 16.
\par
Untuk menyelesaikan soal ini, perhatikan bahwa $Z = (X-10)/3$ adalah versi baku dari
$X$.  Jadi,
\begin{eqnarray*}
P(4 \le X \le 16) & = & P(X \le 16) - P(X \le 4) \\
                  & = & F_X(16) - F_X(4) \\
                  & = & F_Z\left(\frac {16 - 10}3 \right) - F_Z\left(\frac {4-10}3 \right) \\
                  & = & F_Z(2) - F_Z(-2)\ . \\
\end{eqnarray*} 
Ekspresi terakhir ini dapat dihitung menggunakan nilai fungsi distribusi normal baku
dalam tabel (lihat~\ref{app_a}); dengan menggunakan tabel tersebut, kita memperoleh
$F_Z(2) = .9772$ dan $F_Z(-2) = .0228$.  Jadi, jawabannya adalah .9544.
\par
Dalam Bab~\ref{chp 6}, kita akan melihat bahwa parameter $\mu$ adalah rataan, atau
nilai rata-rata, dari peubah acak $X$.  Parameter $\sigma$ mengukur sebaran
peubah acak dan disebut simpangan baku.  Jadi, pertanyaan dalam contoh ini merupakan
pertanyaan khas: berapakah peluang bahwa nilai peubah acak berada dalam jarak dua
simpangan baku dari nilai rata-ratanya?
\end{example}

\subsection*{Densitas Maxwell dan Rayleigh}\index{fungsi densitas!Maxwell}\index{densitas Maxwell}
\index{fungsi densitas!Rayleigh}\index{densitas Rayleigh}

\begin{example}\label{exam 5.19} Misalkan kita menjatuhkan sebuah anak panah pada
permukaan meja yang luas, yang kita anggap sebagai bidang $x$$y$, dan misalkan
koordinat $x$ dan $y$ dari ujung anak panah tersebut saling independen dan
berdistribusi normal dengan parameter $\mu = 0$ dan $\sigma = 1$.  Bagaimanakah
distribusi jarak titik itu dari titik asal?
\par
Persoalan ini muncul dalam fisika ketika diasumsikan bahwa partikel yang bergerak di
$R^n$ memiliki komponen-komponen kecepatan yang saling independen dan berdistribusi
normal, lalu kita ingin mencari densitas laju partikel tersebut.  Densitas untuk kasus
$n = 3$ disebut densitas Maxwell.
\par
Densitas untuk kasus $n = 2$ (yakni percobaan papan sasaran anak panah yang dijelaskan
di atas) disebut densitas Rayleigh.  Kita dapat menyimulasikan kasus ini dengan memilih
secara independen sepasang koordinat $(x,y)$, masing-masing dari distribusi normal
dengan $\mu = 0$ dan $\sigma = 1$ pada $(-\infty,\infty)$, menghitung jarak
$r = \sqrt{x^2 + y^2}$ dari titik $(x,y)$ ke titik asal, mengulangi proses ini
berkali-kali, lalu menyajikan hasilnya dalam grafik batang.  Hasilnya ditampilkan
dalam Gambar~\ref{fig 5.14}.

\putfig{4.5truein}{PSfig5-14}{Distribusi jarak anak panah dalam 1000 jatuhan.}{fig 5.14} 

Kita juga telah memplot densitas teoretis
$$ f(r) = re^{-r^2/2}\ .
$$
Rumus ini akan diturunkan dalam Bab~\ref{chp 7}; lihat Contoh~\ref{exam 7.10}.
\end{example}

\subsection*{Densitas Kai-Kuadrat}\index{densitas kai-kuadrat}\index{fungsi
densitas!kai-kuadrat}

Kita kembali ke persoalan independensi sifat\index{sifat, independensi} yang dibahas
dalam Contoh~\ref{exam 5.6}.  Sering kali kita memiliki dua sifat yang masing-masing
mempunyai beberapa nilai berbeda.  Seperti terlihat dalam contoh tersebut, cukup
banyak perhitungan diperlukan bahkan ketika setiap sifat hanya memiliki dua nilai.
Sekarang kita berikan metode lain untuk menguji independensi sifat yang memerlukan
jauh lebih sedikit perhitungan.
\begin{example}\label{exam 5.20}
Misalkan kita memiliki data dalam Tabel~\ref{table 5.7} mengenai nilai dan gender
mahasiswa dalam sebuah kelas Kalkulus.
\begin{table}
\centering
\begin{tabular}{|c|c|c|c|}
\hline
       &Perempuan\hspace{.15in}&\hspace{.15in}Laki-laki\hspace{.15in}& \\ \hline 
A      & \hspace{.15in}37  & \hspace{.15in}56                &\hspace{.15in}93  \\ \hline 
B      & \hspace{.15in}63  & \hspace{.15in}60                &\hspace{.075in}123\\ \hline 
C      & \hspace{.15in}47  & \hspace{.15in}43                &\hspace{.15in}90  \\ \hline 
Di bawah C& \hspace{.25in}5   & \hspace{.25in}8                 &\hspace{.15in}13  \\ \hline 
       &\hspace{.1in}152   &\hspace{.1in}167                 &\hspace{.075in}319\\ \hline
\end{tabular}
\caption{Data kelas Kalkulus.}
\label{table 5.7}
\end{table}
Dalam situasi ini, kita dapat menggunakan model sejenis dengan yang digunakan dalam
Contoh~\ref{exam 5.6}.  Bayangkan sebuah guci berisi 319 bola dengan dua warna,
misalnya biru dan merah, yang masing-masing mewakili perempuan dan laki-laki.  Kita
ambil 93 bola dari guci tanpa pengembalian.  Bola-bola ini berkaitan dengan nilai A.
Kita lanjutkan dengan mengambil 123 bola yang berkaitan dengan nilai B.  Setelah
selesai, kita mempunyai empat himpunan bola, dengan setiap bola menjadi anggota tepat
satu himpunan.  (Kita juga dapat menetapkan bahwa bola memiliki empat warna yang
berkaitan dengan empat kemungkinan nilai.  Dalam hal ini, kita mengambil himpunan
bagian berukuran 152 yang berkaitan dengan perempuan.  Bola yang tersisa dalam guci
berkaitan dengan laki-laki.  Pilihan ini tidak memengaruhi keputusan akhir apakah
kita harus menolak hipotesis independensi sifat.)
\par
Himpunan data harapan dapat ditentukan dengan cara yang sama persis seperti dalam
Contoh~\ref{exam 5.6}.  Dengan cara ini, kita memperoleh nilai harapan yang
ditampilkan dalam Tabel~\ref{table 5.8}.
\begin{table}
\centering
\begin{tabular}{|c|c|c|c|}
\hline
       & Perempuan \hspace{.15in}&\hspace{.15in}Laki-laki\hspace{.15in}&         \\ \hline 
A      & \hspace{.15in}44.3     & \hspace{.15in}48.7  &\hspace{.15in}93  \\ \hline 
B      & \hspace{.15in}58.6     & \hspace{.15in}64.4  &\hspace{.075in}123\\ \hline 
C      & \hspace{.15in}42.9     & \hspace{.15in}47.1  &\hspace{.15in}90  \\ \hline 
Di bawah C& \hspace{.15in}6.2     & \hspace{.2in}6.8    &\hspace{.15in}13  \\ \hline 
       &  152                   &  167                &\hspace{.075in}319\\ 
\hline
\end{tabular}
\caption{Data harapan.}
\label{table 5.8}
\end{table}
Bahkan jika kedua sifat independen, kita tetap mengharapkan adanya perbedaan di antara
bilangan-bilangan dalam kotak yang bersesuaian pada kedua tabel.  Namun, jika
perbedaannya besar, kita mungkin mencurigai bahwa kedua sifat tersebut tidak
independen.  Dalam Contoh~\ref{exam 5.6}, kita menggunakan distribusi peluang dari
berbagai himpunan data yang mungkin untuk menghitung peluang memperoleh himpunan data
yang menyimpang dari himpunan data harapan sekurang-kurangnya sebesar penyimpangan
data aktual.  Kita dapat melakukan hal yang sama dalam kasus ini, tetapi jumlah
perhitungannya sangat besar.
\par
Sebagai gantinya, kita akan menjelaskan satu bilangan yang dapat mengukur dengan baik
seberapa jauh suatu himpunan data dari himpunan data harapan.  Untuk mengukur jarak
antara kedua himpunan bilangan itu, kita dapat menjumlahkan kuadrat selisih bilangan
yang bersesuaian.  (Kita juga dapat menjumlahkan nilai mutlak selisih, tetapi kita
tidak ingin menjumlahkan selisihnya sendiri.)  Misalkan kita memiliki data yang
memiliki 10 objek harapan dari suatu jenis, tetapi yang kita amati justru 18,
sedangkan dalam kasus lain kita mengharapkan 50 objek dari suatu jenis, tetapi yang
kita amati justru 58.  Meskipun kedua selisih hampir sama, selisih pertama lebih
mengejutkan daripada yang kedua karena banyaknya hasil yang diharapkan dalam kasus
kedua jauh lebih besar daripada dalam kasus pertama.  Salah satu cara mengoreksinya
adalah membagi masing-masing kuadrat selisih dengan bilangan harapan bagi kotak
tersebut.  Jadi, jika nilai dalam delapan kotak pada tabel pertama kita beri label
$O_i$ (untuk nilai teramati) dan nilai dalam delapan kotak pada tabel kedua kita beri
label $E_i$ (untuk nilai harapan), ekspresi berikut cukup masuk akal untuk mengukur
seberapa jauh data teramati dari data harapan:
$$\sum_{i = 1}^8 \frac{(O_i - E_i)^2}{E_i}\ .$$
Ekspresi ini merupakan peubah acak yang biasanya dilambangkan dengan simbol $\chi^2$,
dibaca ``kai-kuadrat."  Nama tersebut digunakan karena, dengan asumsi bahwa kedua
sifat independen, densitas peubah acak ini dapat dihitung dan kira-kira sama dengan
suatu densitas yang disebut densitas kai-kuadrat.  Kita tidak memberikan ekspresi
eksplisit bagi densitas ini karena ekspresi tersebut melibatkan fungsi gamma, yang
belum kita bahas.  Densitas kai-kuadrat sebenarnya merupakan kasus khusus dari
densitas gamma umum.
\par
Dalam penerapan densitas kai-kuadrat, digunakan tabel nilai densitas ini, seperti pada
densitas normal.  Densitas kai-kuadrat memiliki satu parameter $n$, yang disebut
banyaknya derajat kebebasan\index{derajat kebebasan}.  Bilangan $n$ biasanya mudah
ditentukan dari persoalan yang dihadapi.  Sebagai contoh, jika kita menguji
independensi dua sifat, dan kedua sifat masing-masing mempunyai $a$ dan $b$ nilai,
banyaknya derajat kebebasan peubah acak $\chi^2$ adalah $(a-1)(b-1)$.  Jadi, dalam
contoh ini, banyaknya derajat kebebasan adalah 3.
\par
Ingat bahwa dalam contoh ini kita berusaha menguji independensi dua sifat, yaitu
gender dan nilai.  Jika kita mengasumsikan kedua sifat ini independen, model bola dan
guci di atas memberi kita cara untuk menyimulasikan percobaan.  Dengan komputer, kita
telah melakukan 1000 percobaan dan menghitung nilai peubah acak $\chi^2$ untuk setiap
percobaan.  Hasilnya ditampilkan dalam Gambar~\ref{fig 5.14.5}, bersama dengan fungsi
densitas kai-kuadrat dengan tiga derajat kebebasan.
\par
Seperti dinyatakan di atas, jika nilai peubah acak $\chi^2$ besar, kita cenderung
tidak percaya bahwa kedua sifat independen.  Namun, seberapa besar yang disebut
besar?  Nilai aktual peubah acak ini untuk data di atas adalah 4.13.  Dalam
Gambar~\ref{fig 5.14.5}, kita menampilkan densitas kai-kuadrat dengan 3 derajat
kebebasan.  Terlihat bahwa nilai 4.13 lebih besar daripada kebanyakan nilai yang
diambil peubah acak ini.
\putfig{3truein}{PSfig5-14-5}{Densitas kai-kuadrat dengan tiga derajat kebebasan.}{fig 5.14.5}%%4.5truein 

\par
Biasanya, seorang statistikawan menghitung nilai $v$ dari peubah acak $\chi^2$ seperti
yang telah kita lakukan.  Kemudian, dengan melihat tabel nilai densitas kai-kuadrat,
ditentukan suatu nilai $v_0$ yang hanya dilampaui sebanyak 5\% dari seluruh kejadian.
Jika $v \ge v_0$, statistikawan menolak hipotesis bahwa kedua sifat independen.  Dalam
kasus ini, $v_0 = 7.815$, sehingga kita tidak menolak hipotesis bahwa kedua sifat
independen.
\end{example}

\subsection*{Densitas Cauchy}\index{densitas Cauchy}\index{fungsi densitas!Cauchy}
 
Contoh berikut berasal dari Feller.\index{FELLER, W.}\footnote{W. Feller,  \emx {An Introduction
to Probability Theory and Its Applications,}, vol. 2, (New York: Wiley, 1966)}
\begin{example}\label{exam 5.20.5}
Misalkan sebuah cermin dipasang pada sumbu vertikal dan dapat berputar bebas
mengelilingi sumbu tersebut.  Sumbu cermin berjarak 1 kaki dari dinding lurus yang
panjangnya tak berhingga.  Seberkas cahaya dipancarkan ke cermin, lalu sinar
pantulannya mengenai dinding.  Misalkan $\phi$ adalah sudut antara sinar pantulan dan
garis yang tegak lurus terhadap dinding serta melalui sumbu cermin.  Kita asumsikan
bahwa $\phi$ berdistribusi seragam di antara $-\pi/2$ dan $\pi/2$.  Misalkan $X$
menyatakan jarak antara titik pada dinding yang terkena sinar pantulan dan titik pada
dinding yang paling dekat dengan sumbu cermin.  Sekarang kita tentukan densitas $X$.
\par
Misalkan $B$ adalah besaran positif tetap.  Maka $X \ge B$ jika dan hanya jika
$\tan(\phi) \ge B$, yang terjadi jika dan hanya jika $\phi \ge \arctan(B)$.  Hal ini
terjadi dengan peluang
$$\frac{\pi/2 - \arctan(B)}{\pi}\ .$$
Jadi, untuk $B$ positif, fungsi distribusi kumulatif $X$ adalah
$$F(B) = 1 - \frac{\pi/2 - \arctan(B)}{\pi}\ .$$
Dengan demikian, fungsi densitas untuk $B$ positif adalah
$$f(B) = \frac{1}{\pi (1 + B^2)}\ .$$
Karena situasi fisiknya simetris terhadap $\phi = 0$, mudah dilihat bahwa ekspresi
densitas di atas juga benar untuk nilai $B$ negatif.
\par
Hukum Bilangan Besar, yang akan kita bahas dalam Bab~\ref{chp 8}, menyatakan bahwa
dalam banyak kasus, jika kita mengambil rata-rata nilai independen dari suatu peubah
acak, rata-rata itu mendekati bilangan tertentu ketika banyaknya nilai bertambah.
Ternyata, jika hal ini dilakukan pada peubah acak berdistribusi Cauchy, rata-ratanya
tidak mendekati bilangan tertentu mana pun.
\end{example}

\exercises
\begin{LJSItem}

\i\label{exer 5.2.1} Pilih bilangan $U$ dari interval satuan $[0,1]$ dengan
distribusi seragam.  Carilah distribusi kumulatif dan densitas peubah acak
\begin{enumerate}
\item $Y = U + 2$.

\item $Y = U^3$.
\end{enumerate}

\i\label{exer 5.2.2} Pilih bilangan $U$ dari interval $[0,1]$ dengan distribusi
seragam.  Carilah distribusi kumulatif dan densitas peubah acak
\begin{enumerate}
\item $Y = 1/(U + 1)$.

\item $Y = \log(U + 1)$.
\end{enumerate}

\i\label{exer 5.2.2.5} Gunakan Korolari~\ref{cor 5.2} untuk menurunkan ekspresi
peubah acak yang diberikan dalam Persamaan~\ref{eq 5.9}.   \emx {Petunjuk}: Peubah
acak $1 - rnd$ dan $rnd$ berdistribusi identik.


\i\label{exer 5.2.3} Misalkan kita mengetahui peubah acak $Y$ sebagai fungsi dari
peubah acak seragam $U$: $Y = \phi(U)$, dan misalkan kita telah menghitung fungsi
distribusi kumulatif $F_Y(y)$ serta densitas $f_Y(y)$.  Bagaimana kita dapat memeriksa
apakah jawaban kita benar?  Simulasi sederhana memberikan jawabannya: Buatlah grafik
batang $Y = \phi(\mbox{$rnd$})$ dan bandingkan hasilnya dengan grafik $f_Y(y)$.
Kedua grafik seharusnya tampak serupa.  Periksalah jawaban Anda untuk
Latihan~\ref{exer 5.2.1}~dan~\ref{exer 5.2.2} dengan metode ini.

\i\label{exer 5.2.4} Pilih bilangan $U$ dari interval $[0,1]$ dengan distribusi
seragam.  Carilah distribusi kumulatif dan densitas peubah acak
\begin{enumerate}
\item $Y = |U - 1/2|$.

\item $Y = (U - 1/2)^2$.
\end{enumerate}

\i\label{exer 5.2.5} Periksalah hasil Anda untuk Latihan~\ref{exer 5.2.4} melalui
simulasi seperti dijelaskan dalam Latihan~\ref{exer 5.2.3}.

\i\label{exer 5.2.6} Jelaskan cara menghasilkan peubah acak yang fungsi distribusi
kumulatifnya adalah
$$ F(x) =  \left \{ \begin{array}{ll}
              0, & \mbox{jika $x < 0$}, \\
            x^2, & \mbox{jika $0 \leq x \leq 1$}, \\
              1, & \mbox{jika $x > 1.$}
                 \end{array}
        \right.
$$

\i\label{5.2.7} Tulislah program untuk menghasilkan sampel berisi 1000 hasil acak
yang masing-masing dipilih dari distribusi dalam Latihan~\ref{exer 5.2.6}.  Plotlah
grafik batang hasil Anda dan bandingkan densitas empiris ini dengan densitas dari
distribusi kumulatif dalam Latihan~\ref{exer 5.2.6}.

\i\label{exer 5.2.8} Misalkan $U$, $V$ adalah bilangan acak yang dipilih secara
independen dari interval $[0,1]$ dengan distribusi seragam.  Carilah distribusi
kumulatif dan densitas masing-masing peubah
\begin{enumerate}
\item $Y = U + V$.

\item $Y = |U - V|$.
\end{enumerate}

\i\label{exer 5.2.9} Misalkan $U$, $V$ adalah bilangan acak yang dipilih secara
independen dari interval $[0,1]$.  Carilah distribusi kumulatif dan densitas peubah acak
\begin{enumerate}
\item $Y = \max(U,V)$.

\item $Y = \min(U,V)$.
\end{enumerate}

\i\label{exer 5.2.10} Tulislah program untuk menyimulasikan peubah acak dalam
Latihan \ref{exer 5.2.8} dan \ref{exer 5.2.9}, lalu plotlah grafik batang hasilnya.
Bandingkan densitas empiris yang dihasilkan dengan densitas yang ditemukan dalam
Latihan~\ref{exer 5.2.8}~dan~\ref{exer 5.2.9}.

\i\label{exer 5.2.11} Sebuah bilangan $U$ dipilih secara acak dari interval
$[0,1]$.  Carilah peluang bahwa
\begin{enumerate}
\item $R = U^2 < 1/4$.

\item $S = U(1 - U) < 1/4$.

\item $T = U/(1 - U) < 1/4$.
\end{enumerate}

\i\label{exer 5.2.12} Carilah fungsi distribusi kumulatif $F$ dan fungsi densitas $f$
untuk masing-masing peubah acak $R$,~$S$, dan~$T$ dalam Latihan~\ref{exer 5.2.11}.

\i\label{exer 5.2.13} Sebuah titik $P$ dalam persegi satuan memiliki koordinat $X$
dan $Y$ yang dipilih secara acak dari interval $[0,1]$.  Misalkan $D$ adalah jarak
dari $P$ ke sisi persegi terdekat, dan $E$ adalah jarak ke sudut terdekat.  Berapakah
peluang bahwa
\begin{enumerate}
\item $D < 1/4$?

\item $E < 1/4$?
\end{enumerate}

\i\label{exer 5.2.14} Dalam Latihan~\ref{exer 5.2.13}, carilah distribusi kumulatif
$F$ dan densitas $f$ untuk peubah acak $D$.

\i\label{5.2.15} Misalkan $X$ adalah peubah acak dengan fungsi densitas
$$ f_X(x) = \left \{ \begin{array}{ll}
                cx(1 - x), & \mbox{jika $0 < x < 1$}, \\
                        0, & \mbox{selainnya.}
                  \end{array}
         \right.
$$
\begin{enumerate}
\item Berapakah nilai $c$?

\item Apakah fungsi distribusi kumulatif $F_X$ bagi $X$?

\item Berapakah peluang bahwa $X < 1/4$?
\end{enumerate}

\i\label{5.2.16} Misalkan $X$ adalah peubah acak dengan fungsi distribusi kumulatif
$$ F(x) =  \left \{ \begin{array}{ll}
                            0, & \mbox{jika $x < 0$}, \\
              \sin^2(\pi x/2), & \mbox{jika $0 \leq x \leq 1$},  \\
                            1, & \mbox{jika $1 < x$}.
                 \end{array}
        \right.
$$
\begin{enumerate}
\item Apakah fungsi densitas $f_X$ bagi $X$?

\item Berapakah peluang bahwa $X < 1/4$?
\end{enumerate}

\i\label{exer 5.2.17} Misalkan $X$ adalah peubah acak dengan fungsi distribusi
kumulatif $F_X$, serta misalkan $Y = X + b$, $Z = aX$, dan $W = aX + b$, dengan $a$
dan $b$ sebarang konstanta.  Carilah fungsi distribusi kumulatif $F_Y$,~$F_Z$, dan
$F_W$.   \emx {Petunjuk}: Kasus $a > 0$, $a = 0$, dan $a < 0$ memerlukan argumen yang
berbeda.

\i\label{exer 5.2.18} Misalkan $X$ adalah peubah acak dengan fungsi densitas $f_X$,
serta misalkan $Y = X + b$, $Z = aX$, dan $W = aX + b$, dengan $a \ne 0$.  Carilah
fungsi densitas $f_Y$,~$f_Z$, dan~$f_W$.  (Lihat Latihan~\ref{exer 5.2.17}.)

\i\label{exer 5.2.19} Misalkan $X$ adalah peubah acak yang berdistribusi seragam pada
$[c,d]$, dan misalkan $Y = aX + b$.  Untuk pilihan $a$ dan $b$ manakah $Y$
berdistribusi seragam pada $[0,1]$?

\i\label{exer 5.2.20} Misalkan $X$ adalah peubah acak dengan fungsi distribusi
kumulatif $F$ yang naik tegas pada daerah nilai $X$.  Misalkan $Y = F(X)$.  Tunjukkan
bahwa $Y$ berdistribusi seragam pada interval $[0,1]$.  (Rumus $X = F^{-1}(Y)$ lalu
memberi tahu kita cara membentuk $X$ dari peubah acak seragam $Y$.)

\i\label{exer 5.2.21} Misalkan $X$ adalah peubah acak dengan fungsi distribusi
kumulatif $F$.  \emx {Median} $X$ adalah nilai $m$ yang memenuhi $F(m) = 1/2$.  Maka
$X < m$ dengan peluang 1/2 dan $X > m$ dengan peluang 1/2.  Carilah $m$ jika $X$
\begin{enumerate}
\item berdistribusi seragam pada interval $[a,b]$.

\item berdistribusi normal dengan parameter $\mu$ dan $\sigma$.

\item berdistribusi eksponensial dengan parameter $\lambda$.
\end{enumerate}

\i\label{exer 5.2.22} Misalkan $X$ adalah peubah acak dengan fungsi densitas $f_X$.
\emx {Rataan} $X$ adalah nilai $\mu = \int xf_x(x)\,dx$.  Maka $\mu$ memberikan
nilai rata-rata bagi $X$ (lihat Bagian~\ref{sec 6.3}).  Carilah $\mu$ jika $X$
berdistribusi seragam, normal, atau eksponensial, seperti dalam
Latihan~\ref{exer 5.2.21}.

\i\label{exer 5.2.23} Misalkan $X$ adalah peubah acak dengan fungsi densitas $f_X$.
\emx {Modus} $X$ adalah nilai $M$ yang membuat $f(M)$ maksimum.  Maka nilai $X$ di
dekat $M$ paling mungkin terjadi.  Carilah $M$ jika $X$ berdistribusi normal atau
eksponensial, seperti dalam Latihan~\ref{exer 5.2.21}.  Apa yang terjadi jika $X$
berdistribusi seragam?

\i\label{exer 5.2.24} Misalkan $X$ adalah peubah acak berdistribusi normal dengan
parameter $\mu = 70$, $\sigma = 10$.  Estimasikan
\begin{enumerate}
\item $P(X > 50)$.

\item $P(X < 60)$.

\item $P(X > 90)$.

\item $P(60 < X < 80)$.
\end{enumerate}

\i\label{5.2.25} Bridies' Bearing Works memproduksi poros bantalan yang diameternya
berdistribusi normal dengan parameter $\mu = 1$, $\sigma = .002$.  Spesifikasi pembeli
mensyaratkan diameter $1.000 \pm .003$ cm.  Berapa bagian dari poros buatan perusahaan
yang kemungkinan akan ditolak?  Jika produsen memperbaiki kendali mutunya, ia dapat
memperkecil nilai $\sigma$.  Berapakah nilai $\sigma$ yang menjamin bahwa tidak lebih
dari 1~persen porosnya kemungkinan ditolak?

\i\label{exer 5.2.26} Ujian akhir di Podunk University disusun sedemikian rupa
sehingga nilainya kira-kira berdistribusi normal dengan parameter $\mu$ dan $\sigma$.
Pengajar memberikan nilai huruf berdasarkan nilai ujian seperti ditampilkan dalam
Tabel~\ref{table 5.9} (proses ini disebut ``penilaian berdasarkan kurva").
\begin{table}
\centering
\begin{tabular}{lc} Nilai ujian     & Nilai huruf           \\
\hline
$\mu + \sigma < x$                 & A                      \\
$\mu < x < \mu + \sigma$           & B                      \\
$\mu - \sigma < x < \mu $          & C                      \\
$\mu - 2\sigma < x < \mu - \sigma$ & D                      \\
$x < \mu - 2\sigma$                & F                      \\
\end{tabular}
\caption{Penilaian berdasarkan kurva.}
\label{table 5.9}
\end{table}

\noindent Berapa bagian dari kelas yang memperoleh A,~B, C, D,~F?

\i\label{exer 5.2.27} (Ross\footnote{S.~Ross, \emx {A First Course in Probability
Theory,} 2d~ed. (New York: Macmillan, 1984).}) Seorang saksi ahli dalam perkara
paternitas\index{perkara paternitas} bersaksi bahwa lamanya kehamilan (dalam hari),
dari pembuahan hingga persalinan, kira-kira berdistribusi normal dengan parameter
$\mu = 270$, $\sigma = 10$.  Tergugat dalam perkara tersebut dapat membuktikan bahwa
ia berada di luar negeri selama periode dari~290 hingga~240 hari sebelum kelahiran
anak.  Berapakah peluang bahwa tergugat berada di dalam negeri ketika anak itu
dikandung?

\i\label{5.2.28} Misalkan waktu (dalam jam) yang diperlukan untuk memperbaiki mobil
merupakan peubah acak berdistribusi eksponensial dengan parameter $\lambda = 1/2$.
Berapakah peluang bahwa waktu perbaikan melebihi 4 jam?  Jika waktu tersebut melebihi
4 jam, berapakah peluang bahwa waktu itu melebihi 8 jam?

\i\label{5.2.29} Misalkan banyaknya tahun sebuah mobil dapat beroperasi berdistribusi
eksponensial dengan parameter $\mu = 1/4$.  Jika Prosser membeli mobil bekas hari ini,
berapakah peluang bahwa mobil itu masih beroperasi setelah 4 tahun?

\i\label{5.2.30} Misalkan $U$ adalah peubah acak yang berdistribusi seragam pada
$[0,1]$.  Berapakah peluang bahwa persamaan
$$ x^2 + 4Ux + 1 = 0
$$ memiliki dua akar real berbeda $x_1$ dan $x_2$?

\i\label{exer 5.2.31} Tulislah program untuk menyimulasikan peubah acak dengan
densitas berikut, buatlah grafik batang yang sesuai untuk masing-masing densitas, lalu
bandingkan densitas eksaknya dengan grafik batang.

\begin{enumerate}
\item $f_X(x) = e^{-x}\ \  \mbox{pada}\,\, [0,\infty)\,\, 
(\mbox{tetapi\,\,cukup\,\,kerjakan\,\,pada\,\,} [0,10]).$

\item $f_X(x) = 2x\ \ \mbox{pada}\,\, [0,1].$

\item $f_X(x) = 3x^2\ \ \mbox{pada}\,\, [0,1].$

\item $f_X(x) = 4|x - 1/2|\ \ \mbox{pada}\,\, [0,1].$
\end{enumerate}

\i\label{exer 5.2.32} Misalkan kita mengamati suatu proses dengan waktu antar-
kejadian berdistribusi eksponensial dengan $\lambda = 1/30$ (yakni, rata-rata waktu
antar-kejadian adalah 30 menit).  Misalkan proses dimulai pada suatu waktu dan kita
mulai mengamatinya 3 jam kemudian.  Tulislah program untuk menyimulasikan proses ini.
Misalkan $T$ menyatakan lamanya waktu yang harus kita tunggu, setelah mulai mengamati,
sampai sebuah kejadian terjadi.  Buatlah program Anda mencatat $T$.  Berapakah
estimasi nilai rata-rata $T$?

\i\label{exer 5.2.33} Jones memasang dua bohlam baru: bohlam 60~watt dan bohlam
100~watt.  Diklaim bahwa umur bohlam 60~watt memiliki densitas eksponensial dengan
rata-rata umur 200 jam ($\lambda = 1/200$).  Bohlam 100~watt juga memiliki densitas
eksponensial, tetapi rata-rata umurnya hanya 100 jam ($\lambda = 1/100$).  Jones ingin
mengetahui peluang bahwa bohlam 100~watt akan bertahan lebih lama daripada bohlam
60~watt.
\par
Jika $X$ dan $Y$ adalah dua peubah acak independen dengan densitas eksponensial
masing-masing $f(x) = \lambda e^{-\lambda x}$ dan $g(x) = \mu e^{-\mu x}$, maka
peluang bahwa $X$ lebih kecil daripada $Y$ diberikan oleh
$$ P(X < Y) = \int_0^\infty f(x)(1 - G(x))\,dx,
$$ dengan $G(x)$ sebagai fungsi distribusi kumulatif untuk $g(x)$.  Jelaskan mengapa
hal ini berlaku.  Gunakan hasil ini untuk menunjukkan bahwa
$$ P(X < Y) = \frac \lambda{\lambda + \mu}
$$ dan menjawab pertanyaan Jones.

\i\label{exer 5.2.34} Tinjau proses antrean sederhana dalam Contoh~\ref{exam
5.21}.  Misalkan Anda mengamati panjang antrean.  Jika terdapat $j$ orang dalam
antrean, pada perubahan berikutnya panjang antrean akan berkurang menjadi $j - 1$
atau bertambah menjadi $j + 1$.  Gunakan hasil Latihan~\ref{exer 5.2.33} untuk
menunjukkan bahwa peluang panjang antrean berkurang menjadi $j - 1$ adalah
$\mu/(\mu + \lambda)$, sedangkan peluangnya bertambah menjadi $j + 1$ adalah
$\lambda/(\mu + \lambda)$.  Ketika panjang antrean adalah 0, antrean hanya dapat
bertambah menjadi~1.  Tulislah program untuk menyimulasikan panjang antrean.  Gunakan
simulasi ini untuk membantu merumuskan dugaan berupa kondisi pada $\mu$~dan~$\lambda$
yang menjamin bahwa antrean kadang-kadang kosong.

\i\label{exer 5.2.36} Misalkan $X$ adalah peubah acak yang memiliki densitas
eksponensial dengan parameter $\lambda$.  Carilah densitas peubah acak $Y = rX$,
dengan $r$ bilangan real positif.

\i\label{exer 5.2.37} Misalkan $X$ adalah peubah acak yang memiliki densitas normal,
dan tinjau peubah acak $Y = e^X$.  Maka $Y$ memiliki densitas \emx {log-normal}
\index{fungsi densitas!log-normal}\index{densitas log-normal}.  Carilah densitas $Y$ ini.

\i\label{exer 5.2.38} Misalkan $X_1$ dan $X_2$ adalah peubah acak independen, dan
untuk $i = 1, 2$, misalkan $Y_i = \phi_i(X_i)$, dengan $\phi_i$ naik tegas pada
daerah nilai $X_i$.  Tunjukkan bahwa $Y_1$ dan $Y_2$ independen.  Perhatikan bahwa
hasil yang sama berlaku tanpa asumsi bahwa $\phi_i$ naik tegas, tetapi pembuktiannya
lebih sulit.

\end{LJSItem}
%\end{LJSItem}}
%\end{document}
